\chapter{引言}

\section{研究背景与意义}

\subsection{大语言模型的崛起与能力演进}

近十年间，人工智能研究在深度学习推动下持续演进，其中大语言模型（Large Language Models，LLMs）取得了快速发展。从早期的 BERT、GPT-2 到近年来的 GPT-4、Llama 系列、Qwen 系列和 DeepSeek，模型参数规模由亿级扩展至千亿级，训练语料由数百 GB 扩展至数十 TB，模型在语言理解、知识表征与文本生成等方面的能力随之提升。

这场技术演进的起点，是 Vaswani 等人于 2017 年提出的 Transformer 架构。Transformer 以自注意力机制（Self-Attention）取代了循环神经网络对序列信息的逐步传递，使得模型在捕捉长程依赖关系上具有了结构性的先天优势。在此基础上，以"预测下一个 Token"为核心任务的自回归语言模型（Autoregressive Language Model）确立了大规模预训练的基本范式：通过在海量互联网文本上最小化负对数似然损失，模型被迫学习并压缩了关于人类语言、世界知识乃至逻辑结构的大量信息。

然而，仅仅通过预训练阶段获得的模型，在实际应用中存在明显的局限性。预训练模型的优化目标是对自然语言分布建模，而非针对特定任务的准确完成。这导致原始的预训练模型在执行用户指令时往往缺乏针对性，甚至会生成偏离预期、内容有害的输出。为了解决这一问题，监督微调（Supervised Fine-Tuning，SFT）技术被引入到 LLM 的训练流程之中。

监督微调的核心思路是：在人工标注的高质量"问题—答案"对数据上，以最大似然估计的方式对模型进行进一步训练，使其学会"模仿"专家给出的示范性回答。该方法能够提升模型的指令遵循能力，并使输出更符合人类阅读预期。

但是，SFT 方法的性能仍然受训练数据质量与覆盖范围的显著约束。对于数学推导、逻辑推理、代码生成等高难度任务，构建规模充足且质量稳定的专家标注数据代价较高。更重要的是，SFT 的优化目标仍然主要对应于示范分布拟合，而非显式的推理过程优化。因此，模型往往更擅长复现训练数据中已有的解答模式，而在分布外变形问题上容易出现性能下降。这一现象在数学竞赛题和算法题等强逻辑任务中尤为明显。

\subsection{强化学习后训练的兴起}

在此背景下，强化学习（Reinforcement Learning，RL）被逐步引入大语言模型的后训练阶段，并成为提升模型复杂推理能力的重要技术路径。

强化学习后训练（Reinforcement Learning Fine-Tuning，RLFT）的基本思路与监督微调有着本质的不同。SFT 要求模型"模仿答案"，而 RL 则要求模型通过与环境的交互来"探索策略"。在 LLM 的语境下，这个"环境"可以是人类评估者给出的偏好反馈，也可以是一个自动化的验证器（Verifier）。对于数学、代码等具有客观正确答案的任务，验证器可以由规则系统直接构成：答案是否正确，程序是否通过测试，逻辑链条是否自洽。这类基于可验证奖励的强化学习方法，被称为 RLVR（Reinforcement Learning with Verifiable Rewards），近年来获得了学术界和工业界的广泛关注。

在强化学习的框架中，LLM 被视为一个策略（Policy），其参数 θ 决定了在给定输入 Prompt x 下输出回答 a 的概率分布 π_θ(a|x)。强化学习的目标是最大化策略在所有输入上的期望奖励：

\begin{equation}
\label{eq:rlvr_objective}
J(\theta) = \mathbb{E}_{x \sim d_0,\; a \sim \pi_{\theta}(\cdot\mid x)}\bigl[r(x,a)\bigr]
\end{equation}

其中 d_0 是 Prompt 的采样分布，r(x, a) 是验证器给出的奖励信号。通过策略梯度（Policy Gradient）方法，我们可以在不知道奖励函数梯度的情况下，利用采样得到的回答对模型参数进行迭代优化。

这一范式的重要性在于，模型不再完全受限于人类专家示范数据，而可以在输出空间中基于奖励信号进行策略搜索。只要验证器能够提供稳定且可判别的奖励，模型便有可能提升高奖励推理路径的生成概率。因此，强化学习后训练已成为近年来提升前沿模型推理能力的重要组成部分。

\subsection{规模化训练中的计算代价}

然而，强化学习后训练通常伴随显著更高的计算成本。与监督微调相比，RLFT 的额外开销主要来自在线采样（On-policy Sampling）的必要性。

在策略梯度方法中，模型必须基于当前策略 π_θ 持续生成新的样本，用于估计梯度方向。这是因为策略梯度定理要求所使用的样本必须来自于当前策略的分布——若使用旧策略生成的样本，梯度估计将产生严重的偏差，导致训练不稳定甚至发散。这就意味着，每一轮梯度更新之前，都必须进行一轮代价高昂的推理过程（即让模型生成文本），才能获得用于计算梯度的原始数据。

在实际测量中，对于主流的数学推理任务，单轮训练步骤中推理采样所消耗的计算量通常占整体训练算力的 80% 以上。这意味着，训练效率的瓶颈并不在于参数的反向传播，而在于如何高效地从当前策略中获取有价值的样本。

更进一步，主流算法如 GRPO（Group Relative Policy Optimization）为了估计每个 Prompt 的相对优势，通常对每个输入生成 n 个候选回答（n 一般设置为 8 至 32 之间），通过组内相对奖励来计算优势函数。这种"以采样换精度"的策略虽然能够降低方差、稳定训练，却同时也将推理计算的开销再次放大了 n 倍。

因此，如何在保证训练信号质量的前提下，提高单位计算资源对应的有效梯度贡献，已成为 LLM 强化学习后训练中的关键问题。


\section{核心问题：信号坍塌与采样低效}

\subsection{GRPO 框架与信号坍塌现象}

在 GRPO 框架下，对于一个给定的 Prompt $x$，模型生成 $n$ 个回答 $\{a_1,a_2,\ldots,a_n\}$，验证器给出对应的奖励向量 $\mathbf{r}=[r_1,r_2,\ldots,r_n]$。每个回答的优势函数可定义为：
\begin{equation}
\label{eq:grpo_adv}
A_{\mathrm{GRPO}}(x,a_i) = \frac{r_i-\bar r}{\sigma_r+\varepsilon},\qquad
\bar r = \frac{1}{n}\sum_{j=1}^{n} r_j,\quad
\sigma_r = \sqrt{\frac{1}{n}\sum_{j=1}^{n}(r_j-\bar r)^2}.
\end{equation}
其中 $\varepsilon$ 为防止除零的小量。

由式\eqref{eq:grpo_adv} 可知，当 $n$ 个回答全部正确（$r_i=1,\forall i$）或全部错误（$r_i=0,\forall i$）时，组内奖励标准差 $\sigma_r=0$ 且 $r_i=\bar r$，从而 $A_{\mathrm{GRPO}}(x,a_i)=0$，梯度为零，模型参数不发生更新。本文将该现象称为信号坍塌（Signal Collapse）。

信号坍塌并非罕见的边界情形，而是固定采样预算下普遍存在的结构性问题。对于较易题目，随着训练推进，模型生成正确回答的概率逐渐接近 1，此时固定 n 次采样更容易全部正确并产生正向坍塌。对于较难题目，当单次采样命中正确答案的概率较低时，有限预算下更容易出现全部错误并产生负向坍塌。两类情形都会使对应 Prompt 在当前训练步中无法提供有效更新信号。

实验数据表明，该问题在训练过程中具有较高发生比例。如表 1.1 所示，在以 Qwen2.5-Math-1.5B 为基础模型、使用 GRPO 在数学推理数据集上训练的实验中，我们统计了不同采样规模 n 下，训练前期与后期中发生信号坍塌的 Prompt 比例。

\begin{table}[htbp]
\centering
\caption{不同采样规模下的信号坍塌比例（示例统计）}
\label{tab:collapse_ratio}
\small
\begin{tabular}{ccc}
\hline
采样规模 $n$ & 训练初期坍塌率(\%) & 训练后期坍塌率(\%) \\
\hline
4  & 42.1 & 78.5 \\
8  & 28.3 & 61.2 \\
16 & 15.6 & 54.8 \\
32 & 8.2  & 50.4 \\
\hline
\end{tabular}
\end{table}

表 1.1 表明，即便将采样规模提升至 n = 32，训练后期仍有超过一半的 Prompt 无法提供有效更新信号。与此同时，将 n 从 4 提升到 32 会使推理计算开销增加 8 倍，而坍塌率的下降幅度并不与成本增长同比例。这说明，仅通过增加固定采样数量缓解信号坍塌，其计算效率是有限的。

\subsection{Pass@k 实验：上限被遮蔽，而非不存在}

信号坍塌现象进一步引出一个问题：当某个 Prompt 在训练中频繁发生负向坍塌时，这究竟意味着模型不具备解题能力，还是仅仅意味着有限采样预算不足以显式观测到低概率正确路径？

Pass@k 实验为这一问题提供了经验依据。在基于 Qwen2.5-Math-1.5B 的一组系统性实验中，我们固定模型权重（使用训练中期的 checkpoint），分别统计不同采样次数 k 下模型在数学测试集上至少生成一个正确回答的题目比例。结果表明：

当 k = 1 时，Pass@k 仅为 26.5%，意味着大多数题目在单次采样下模型无法给出正确答案。
而当 k = 256 时，Pass@256 达到了 81.3%，说明绝大多数题目在充分采样的条件下，模型是具备解题能力的。

上述差异说明，在不少情形下，问题并不在于模型完全缺乏解题能力，而在于正确路径虽然已存在于策略分布之中，却因概率质量较低而难以在小样本采样中被观测到。

此外，我们进一步比较了不同训练阶段 checkpoint 在 k = 4 与 k = 256 下的表现差异。在训练较早的 checkpoint 上，n = 4 时有 35.3% 的 Prompt 触发正向坍塌（全部正确），而在 n = 256 时该比例降至 10.2%。这表明，一部分在小采样预算下表现稳定的题目，在更大采样规模下仍可能存在尚未充分覆盖的困难情形。

这些观察与已有 pass@k 分析结论具有一致性。例如，Yang Yue 等人在 2025 年的研究（arXiv:2504.13837）指出，RLVR 的主要作用更接近于在已有分布内放大高奖励路径的概率质量，而非在基座模型能力边界之外生成全新的推理路径；因此，模型在 pass@1 上的提升未必对应充分采样条件下能力上界的同步扩展。

本文认为，上述现象的关键约束之一在于主流训练范式中的采样效率不足。在标准 GRPO 训练中，固定且有限的采样预算使大量低概率但真实存在的正确路径长期处于统计上难以观测的状态，模型因而无法稳定获得与这些路径相关的梯度信号，也难以持续提升其概率质量。这一观察构成了本文方法设计的直接动机。

\subsection{负向坍塌概率与有效更新频率分析}

为进一步量化固定采样预算对难题学习效率的影响，下面对负向坍塌概率及其对应的有效更新频率进行分析。

设某 Prompt $x$ 在当前策略下的内生正确率为 $p$，即从 $\pi_{\theta}(\cdot\mid x)$ 中独立采样的每个回答以概率 $p$ 为正确答案。假设每轮对该 Prompt 采样 $n$ 次，则 $n$ 次采样全部错误（负向坍塌）的概率为：
\begin{equation}
\label{eq:collapse_prob}
P_{\mathrm{collapse}}=\Pr(K=0)=(1-p)^n.
\end{equation}

定义有效训练间隔 $T$ 为平均需要多少轮训练，该 Prompt 才能获得一次非零梯度更新机会，则
\begin{equation}
\label{eq:effective_interval}
\mathbb{E}[T]=\frac{1}{1-P_{\mathrm{collapse}}}=\frac{1}{1-(1-p)^n}.
\end{equation}

当 $pn\ll 1$ 时，利用一阶泰勒展开 $(1-p)^n\approx 1-np$，可得近似
\begin{equation}
\label{eq:effective_interval_approx}
\mathbb{E}[T]\approx \frac{1}{np}.
\end{equation}

当 $n\ll 1/p$ 时，$\mathbb{E}[T]$ 将随 $n$ 的不足而快速增大。以一道内生正确率 $p=1/256\approx 0.0039$ 的题目为例，若采样数 $n=8$，则
\begin{equation}
P_{\mathrm{collapse}}=\left(1-\frac{1}{256}\right)^8\approx 1-\frac{8}{256}=96.875\%,\qquad
\mathbb{E}[T]\approx \frac{1}{8\times (1/256)}=32.
\end{equation}

这意味着，在总训练步数为 600 步的训练计划中，该题平均仅能获得 600 / 32 ≈ 19 次非零梯度更新机会。若奖励信号本身还存在较强随机性，这样的有效更新次数通常不足以支撑稳定的参数修正。

进一步地，增大固定采样数 n 对高难题的改善受到真实正确率 p 的显著限制。将 n 从 8 提升到 16 时，若 p = 1/256，则 E[T] 可由 32 降至 16；但若 p = 1/1024，则在 n = 16 时仍有 E[T] ≈ 64。换言之，当 p 足够小时，仅依赖固定预算扩张并不能有效缩短难题的学习间隔，而对应计算成本却近似线性增长。

上述分析表明，缓解难题上的信号坍塌问题，关键不在于对所有样本等比例放大 n，而在于依据样本难度动态分配采样资源。


\section{现有解决方案及其不足}

面对信号坍塌问题，现有研究提出了若干解决方案，但每种方案都存在不可忽视的局限性。

被动过滤（Passive Filtering）是较为直接的应对策略：在每轮训练中，若某个 Prompt 的采样结果发生坍塌，则直接将其从当前 batch 中剔除，仅利用产生有效信号的 Prompt 计算梯度。该方法虽然避免了零梯度样本直接进入更新，但会引入训练分布偏移。随着模型能力提升，简单题更容易因正向坍塌被剔除，难题则更容易因负向坍塌被剔除，最终保留下来的往往是中等难度样本。此外，被动过滤并不减少已发生的推理成本。

固定大规模采样（Large Fixed-Budget Sampling）是另一类常见思路，即通过将采样数 n 增大至 64、128 乃至 512 来降低坍塌概率。该方法在理论上可以逼近充分采样的极限，但计算代价极高。在 8 张 H100 GPU 的训练配置下，将 n 从 8 提升至 128 会使单步训练时间由约 3 分钟增至接近 50 分钟，因此难以作为常规训练方案。

两阶段预算分配（Two-stage Budget Allocation）是一种更精细化的方案：先使用少量样本估计每个 Prompt 的通过率，再据此分配后续采样预算，将更多资源分配给通过率适中的题目。然而，这类方法同样存在局限：估计阶段本身需要额外计算资源，小样本条件下的通过率估计也不稳定；此外，两阶段串行流程与现代推理引擎的流水线并行架构并不完全兼容，工程实现复杂度较高。

已有自适应采样工作提出，不再为所有 Prompt 预先设定统一且固定的采样数 n，而是分轮追加采样，并在满足预设退出条件时提前停止。以代表性工作 Reinforce-Ada（arXiv:2510.04996）为例，该类方法通常在每轮采样固定数量样本（如 M = 4）后，判断是否满足有效信号条件，例如 Balance 条件，即至少同时观测到一个正样本和一个负样本。此类机制能够减少简单题上的冗余采样，并为难题保留进一步追加预算的可能。

然而，现有自适应采样方法在以下三个方面仍存在进一步改进空间，这也是本文关注的核心问题：

第一，Balance 条件在训练后期容易对简单题引入额外开销。当模型对某些题目的通过率已接近 1 时，为满足同时观测正负样本的停止条件，往往需要继续采样以寻找少量负样本，从而增加不必要的计算成本。

第二，固定小批次追加采样与现代推理引擎的 KV Cache 并行架构之间存在适配问题。若每轮仅采样少量样本并立即判断是否退出，则需要频繁重新调度 Prefill 计算，前缀复用率受限，推理吞吐量难以充分释放。

第三，当采样轮次增加时，适度提高采样温度有助于提升对低概率正确路径的探索能力，但温度变化会引入采样分布与策略分布之间的偏差。若不进行校正，梯度估计可能产生系统性偏差。

总体而言，现有方案尚未同时兼顾采样效率、硬件友好性与梯度估计正确性。本文旨在围绕这三个维度提出一套统一的改进框架。


\section{研究内容与主要贡献}

针对上述问题，本文开展了以下三个层面的研究工作：

在理论层面，本文从信息论视角建立了题目难度与单位采样成本信息效率之间的定量关系，证明了难度感知的预算分配相较于难度无关的均匀分配具有更高的信息增益效率。基于此，本文提出以"正信号优先（Positive-first）"退出条件替代"平衡条件（Balance）"作为训练后期的采样停止准则，并给出相应的理论分析与实验验证。

在系统层面，本文提出了指数增长批次采样策略（Exponential Batch Sampling，EBS），令累计采样预算按轮次呈几何扩展；在本文默认设置下，对应的逐轮批次为 2、2、4、8、16，即初始两轮保持相同规模，此后按 2 倍增长。该设计在保留自适应退出灵活性的同时，提高了推理引擎中的 KV Cache 前缀复用率，从而在固定算力预算下缩短单步训练时间。

在算法层面，本文提出层级几何自适应采样框架（Hierarchical Geometric Adaptive Sampling，H-GAS）中的温度分层重要性采样（Temperature-Stratified Importance Sampling，TSIS）机制，将动态温度缩放（Dynamic Temperature Scaling）与重要性采样校正统一到同一套策略梯度框架之中。具体而言，本文令采样温度随累计采样深度对数增长，以提升深度采样阶段对低概率正确路径的探索能力；同时，对来自不同温度的样本统一引入重要性权重，对策略梯度项与基准线估计进行联合校正。该机制的关键意义在于：重要性权重不仅用于修正多温度采样带来的分布偏差，也隐式刻画了样本相对于目标策略的"可信度"与"训练价值"。因此，低温下高置信生成但最终错误的样本会获得更强的惩罚信号，而高温下偶然探索到的正确样本则会被适度折现，避免训练过程过度偏向侥幸解。

上述三项贡献共同构成层级几何自适应采样框架 H-GAS：其中，Positive-first 对应解决第一章 1.3 中的"Balance 条件后期额外开销"问题，EBS 对应解决"固定小批次与 KV Cache 不适配"问题，TSIS 对应解决"温度变化引入分布偏差未校正"问题。整体框架兼顾采样效率、系统吞吐与统计正确性，并在 Qwen2.5-Math-1.5B、Qwen2.5-Math-7B 等开源数学推理模型上进行了实验评估。


\section{论文组织结构}

本论文共六章，各章内容安排如下：

第一章为引言，介绍研究背景、核心问题、现有方法局限与本文贡献。

第二章为相关工作综述，梳理大语言模型强化学习后训练的技术演进、采样效率优化方向的现有研究，以及推理引擎并行加速技术的相关工作。

第三章为方法论，系统阐述本文提出的渐进式自适应采样框架，包括难度感知的信息效率分析、指数增长批次策略的设计原理，以及动态温度缩放与重要性采样校正的算法推导。

第四章为实验，介绍实验设置，并分别从主要性能对比、模块消融分析、计算效率验证及训练动态分析四个维度呈现实验结果。

第五章为讨论，深入分析各方法模块背后的作用机制，并指出本文研究的局限性与未来改进方向。

第六章为结论，总结全文并展望后续研究。


\cleardoublepage
\chapter{相关工作}

\section{大语言模型强化学习后训练与可验证奖励范式}

\subsection{从监督微调到 RLHF 再到 RLVR}

大语言模型的训练通常经历预训练、监督微调与后训练三个阶段。其中，预训练负责学习一般性的语言统计规律与知识表示，监督微调进一步塑造模型的指令跟随能力，而后训练则通过偏好信号、奖励信号或环境反馈，对模型行为进行面向目标场景的再优化。[待补文献：LLM训练三阶段综述] 在复杂推理任务中，尤其是在数学、代码与形式逻辑等“结果可验证”场景中，强化学习后训练（Reinforcement Learning for Fine-Tuning/Post-Training, RLFT）已经成为提升模型推理能力的重要技术路线。[待补文献：RLFT/RLVR综述]

在这一演进链条中，监督微调（SFT）负责把预训练模型从通用语言建模任务迁移到指令跟随场景，但其根本机制仍然是对高质量示范分布的模仿，因此对数据质量与覆盖范围高度敏感。对于数学、代码与复杂推理任务而言，SFT 面临两个突出问题：一是高质量过程监督数据昂贵且稀缺，二是即便拥有充足示范，模型学到的往往仍是“像训练集一样回答”，而不是稳定可迁移的推理策略。[待补文献：InstructGPT；待补文献：Alpaca/Vicuna/Llama-Chat]

为了突破纯监督模仿的局限，RLHF 进一步将后训练形式化为基于人类偏好的策略优化问题。其经典流程通常先以 SFT 初始化策略模型，再用人类比较数据训练奖励模型，最后通过 PPO 等策略优化算法提升被偏好输出的概率。[待补文献：InstructGPT；待补文献：PPO in RLHF] RLHF 在通用对话、安全性与帮助性等维度取得了显著成功，但其上限又受制于奖励模型本身的质量，容易受到偏好噪声、数据规模限制与 reward hacking 的影响。

针对这些问题，基于可验证奖励的强化学习（RLVR）提供了更适合推理任务的替代路径。RLVR 不再依赖主观偏好模型，而是直接调用可执行验证器、答案比对器、单元测试或符号规则，对模型输出进行自动判分。这一路线的代表性成果表明，仅依赖结果可验证奖励，也能在数学推理等任务上显著激发模型潜在能力。[待补文献：DeepSeek-R1-Zero；待补文献：Verifier-based RL]

与传统监督学习不同，强化学习后训练不要求模型给出唯一标准的中间推理链条，而是允许模型围绕同一输入生成多个候选回答，并利用外部验证器对结果进行自动打分。对于数学任务，这一验证器通常由答案比对规则、程序执行器或符号计算器构成；对于代码任务，则可由单元测试、执行结果或形式约束进行评估。[待补文献：Verifier-based RL] 这一范式通常被称为可验证奖励强化学习（Reinforcement Learning with Verifiable Rewards, RLVR），其优势在于无需人工逐条构造细粒度监督信号，能够直接围绕最终任务目标进行优化。

在优化方法上，当前大语言模型后训练广泛采用基于策略梯度的近端更新框架，例如 PPO、RLOO 及其针对大语言模型场景的变体。[待补文献：PPO；待补文献：RLOO；待补文献：GRPO] 其中，Group Relative Policy Optimization（GRPO）由于无需显式价值网络、实现简洁、与结果可验证任务适配性较强，而成为当前推理类后训练中的代表性方法之一。GRPO 对每个 Prompt 采样固定数量的回答，并根据组内奖励差异构造相对优势。该思路在工程上具有较高可扩展性，但也把训练信号是否存在这一问题，直接绑定到了有限样本组内是否能观察到奖励差异之上。

\subsection{关于强化学习能力边界的争论}

围绕 RL 是否真正“创造”了新的推理能力，近期也出现了具有代表性的外部讨论。\texttt{Does Reinforcement Learning Really Incentivize Reasoning Capacity in LLMs Beyond the Base Model?} 系统比较了多类 RLVR 算法与对应基座模型在大 \(k\) 条件下的 Pass@k 表现，并指出当前 RLVR 更多是在低采样预算下提升了正确推理路径被触发的概率，而非显著突破基座模型原有能力边界。[待补文献：Does Reinforcement Learning Really Incentivize Reasoning Capacity in LLMs Beyond the Base Model?] 这一结论虽然并不否定 RLFT 的价值，却对“采样是否足够充分”提出了更尖锐的问题：如果 RL 的主要作用是激发基座模型已具备但未充分显露的能力，那么采样效率就不再只是训练加速问题，而会直接决定这些潜在能力是否能被训练过程观察到、放大并稳定固化。

围绕这一论点，后续研究又出现了两类重要跟进。其一，\texttt{Breaking the Capability Ceiling of LLM Post-Training by Reintroducing Markov States} 明确承认现有 post-training 可能面临能力天花板，并尝试从训练状态建模层面重新打开能力扩张空间。[待补文献：Breaking the Capability Ceiling of LLM Post-Training by Reintroducing Markov States] 其二，\texttt{The Debate on RLVR Reasoning Capability Boundary: Shrinkage, Expansion, or Both? A Two-stage Dynamic View} 则试图调和“边界收缩”与“边界扩张”两种看法，提出更具阶段性的动态解释框架。[待补文献：The Debate on RLVR Reasoning Capability Boundary: Shrinkage, Expansion, or Both? A Two-stage Dynamic View] 这些研究的共同意义在于，它们都把注意力重新引向一个被长期低估的问题：训练过程能否真正观察到模型已有但低概率出现的正确路径。

正是在这一背景下，采样机制从原本的“数据获取步骤”上升为影响训练效率、稳定性与性能上限的核心组成部分。换言之，在 RLFT 中，优化器是否能够学习到有效策略，不仅取决于奖励定义是否合理，还取决于采样过程是否真正触达了模型已存在但尚未稳定外显的行为差异。这一认识构成了本文相关工作的共同出发点。

\section{针对无效信号的早期解决思路}

围绕组相对优化中的无效样本与零梯度问题，现有工作最早形成了三类较为直接的应对思路：被动过滤、暴力大采样以及两阶段预算分配。

第一类方法是\textbf{被动过滤}。这类方法在完成 rollout 之后，直接删除那些组内奖励完全一致、无法构造有效优势的 Prompt，仅利用剩余样本更新模型；\texttt{DAPO} 也可归入这一思路，因为它同样会在组内奖励无差异时跳过对应 Prompt 的更新。[待补文献：Passive Filtering/RAFT/Reinforce-Rej/DAPO相关] 此类思路的优点是实现极其简单，且能够减少零梯度样本对训练统计量的干扰。然而，其局限也十分明显：一方面，推理成本已经实际发生，过滤只能减轻更新噪声，无法回收前向采样开销；另一方面，对于在小样本下长期呈现“全错”的困难题而言，被动过滤等价于反复将最需要学习的样本排除在优化过程之外，导致训练分布向“容易产生信号”的样本倾斜。

第二类方法是\textbf{暴力大采样}。该思路不改变训练框架本身，而是通过显著增大组大小，提升同一 Prompt 下同时采到正确与错误样本的概率，从而缓解零标准差问题。[待补文献：Large-rollout GRPO相关实践] 这一方案的逻辑直观且通常有效，但其代价是推理成本、显存占用与系统调度压力急剧上升。对于参数规模较大的推理模型而言，单纯将每个 Prompt 的 rollout 数提升至数百甚至更高，往往难以作为常规训练配置长期维持。因此，这类方法更像是“用算力换信号”，而不是从机制上提升单位成本的信息获取效率。

第三类方法是\textbf{两阶段预算分配}。其基本思想是先用一部分预算估计样本难度、通过率或信息价值，再依据估计结果为不同 Prompt 分配后续 rollout 预算。[待补文献：Two-stage allocation / GVM-RAFT / 相关预算分配工作] 从统计设计角度看，这种方法已经开始意识到“不同样本不应被一视同仁”；但在大语言模型后训练场景中，这类方法往往仍面临三个困难：其一，前置估计本身需要额外样本，存在前期预算浪费；其二，小样本估计误差较大，预算分配容易受到噪声误导；其三，离线估计或准离线估计流程不易与大规模在线训练系统紧密耦合。

总体而言，这三类早期思路已经揭示了一个基本事实：固定 rollout 数并非最优设计，采样预算需要围绕“哪里更可能产生学习信号”进行重新组织。但它们仍然要么停留在采样之后的被动补救，要么依赖高成本的大样本估计，尚未充分解决在线训练场景下的信号效率问题。

\subsection{前置选择性过滤：在 rollout 之前跳过低价值 Prompt}

在上述“采样后再处理”的思路之外，近期研究进一步提出，应尽可能在 rollout 开始之前就识别并跳过低信息量样本。\texttt{Act Only When It Pays: Efficient Reinforcement Learning for LLM Reasoning via Selective Rollouts} 提出的 \texttt{GRESO} 正是这一方向的代表工作。[待补文献：GRESO] 该工作观察到，零方差或近零方差 Prompt 在训练过程中具有较强的时间一致性：如果某个 Prompt 在当前阶段很难提供有效信号，那么它在相邻训练阶段中大概率仍然缺乏训练价值。基于这一发现，GRESO 利用历史 reward dynamics 构造轻量级的 pre-rollout filtering 机制，在正式生成回答前预测哪些 Prompt 可能是“无信息量”的，并对其进行概率性跳过。

与被动过滤相比，GRESO 的关键区别在于它把“是否值得 rollout”前移为显式决策，从而真正回收了生成开销而不仅仅是更新开销。该工作报告，在多个数学推理模型与基准上，GRESO 可在无明显精度损失的前提下获得最高约 \texttt{2.4×} 的 rollout wall-clock 加速与约 \texttt{2.0×} 的整体训练加速。这一结果说明，围绕零信号样本的优化并不一定都要通过增大采样数实现；在某些场景下，直接减少对低价值 Prompt 的前向投入反而更有效。与此同时，这一路线也存在其局限：它更擅长识别“当前看起来不值得采”的样本，却未必能细致刻画不同难度样本之间的相对信息收益，更难与后续的批次组织与温度调度形成统一设计。

\section{在线自适应采样：从固定组大小到逐轮追加预算}

为了克服被动过滤和两阶段预算分配的局限，近期研究开始转向\textbf{在线自适应采样}。这一方向的核心思想是：不再在训练开始前为每个 Prompt 固定分配相同数量的回答，而是在采样过程中根据已观测到的奖励结构，动态决定是否继续为该 Prompt 追加 rollout。当系统认为某个 Prompt 已经暴露出足够训练信号时提前停止，而当信号仍不充分时则继续投入预算。该思路直接把“是否继续采样”视为训练过程中的在线决策问题。

\texttt{Reinforce-Ada: An Adaptive Sampling Framework under Non-linear RL Objectives} 是这一方向中极具代表性的工作。[待补文献：Reinforce-Ada] 该方法不采用传统两阶段“先估计后分配”的设计，而是通过 successive elimination 式的在线过程，将估计与采样交织进行。它在每一轮仅对仍处于活跃状态的 Prompt 继续采样，并在达到特定退出条件后停止对该 Prompt 的进一步预算投入。其重要贡献在于明确指出：GRPO 等固定组大小方法中的信号缺失并不一定意味着 Prompt 本身没有可学习信息，很多时候只是小样本下未能观测到正负奖励差异而已。通过逐轮追加预算，自适应采样能够在同等总体开销附近恢复这部分被统计假象掩盖的学习信号。

\texttt{Reinforce-Ada} 的另一个启发在于，它把研究重点从“如何修改优势函数”转向“如何更高效地让优势函数非零”。这一转向具有方法论意义：与其在信号已经缺失之后设计复杂的补偿项，不如在采样阶段尽量避免让信号消失。也正因如此，在线自适应采样已逐渐成为结果可验证 RLFT 中提升训练效率的关键方向。

然而，现有在线自适应采样研究也存在若干未被充分解决的问题。首先，许多方法仍默认不同难度 Prompt 应遵循相同的退出逻辑，即只要满足某种平衡条件就认为采样足够，而没有进一步追问：在相同计算成本下，简单题与难题的单位信息收益是否真的相同。其次，多轮采样在统计上虽然更高效，但其常见实现方式往往忽略了现代推理引擎中的缓存复用特性。再次，当采样规模在不同轮次间动态变化时，不同轮次样本的探索强度是否需要同步调整，以及由此带来的分布偏差如何校正，也尚未形成系统解法。本文的工作正是在这些空白之上继续展开。

\section{难度感知与能力对齐：围绕“哪些 Prompt 值得追加预算”的研究}

在自适应采样方向进一步发展后，研究者开始意识到，仅仅做到“有信号就继续、没信号就停止”仍然不够，因为不同 Prompt 的信息密度并不相同。由此，一批工作转向更细粒度地研究\textbf{难度感知}、\textbf{能力对齐}与\textbf{数据中心化预算分配}问题。

\subsection{基于难度的动态 rollout 分配}

\texttt{Enhancing Efficiency and Exploration in Reinforcement Learning for LLMs} 从难度感知 rollout 分配与探索保持两个角度切入，指出现有 RL 方法为所有问题分配相同回答数并不高效。[待补文献：Enhancing Efficiency and Exploration in Reinforcement Learning for LLMs] 该工作认为，简单问题继续投入额外 rollout 通常收益有限，而困难问题则需要更多样本才能暴露出潜在正确路径。因此，它依据题目难度动态调整 rollout budget，并同时引入动态温度调节与退火机制，以维持训练过程中的策略熵，避免 RL 过早丧失探索能力。该工作的重要意义在于：它不再把采样效率与探索能力视为彼此冲突，而是强调二者需要联动设计。

\texttt{Depth-Breadth Synergy in RLVR: Unlocking LLM Reasoning Gains with Adaptive Exploration} 则提出了更具针对性的 \texttt{Difficulty Adaptive Rollout Sampling (DARS)}。[待补文献：DARS] 该工作指出，GRPO 式组相对更新在统计上会系统性低估低成功率、高难度样本的重要性，导致训练过于偏向中等准确率问题，从而限制高难度推理能力的释放。为了解决这一偏置，DARS 通过多阶段 rollout 对困难样本进行重加权与追加采样，使得高难度问题能够获得更多正样本暴露机会。更进一步，该工作把“depth”与“breadth”作为两个正交维度分开讨论：前者对应单题采样深度，主要影响 Pass@K；后者对应训练中覆盖的问题广度，主要影响 Pass@1。这一分析把采样问题从单纯的 rollout 数选择，扩展为探索深度与训练覆盖面的联合设计问题。

与上述方法相比，\texttt{Adaptive Rollout Allocation for Online Reinforcement Learning with Verifiable Rewards} 更进一步从梯度方差最小化角度给出 rollout 分配原则。[待补文献：Adaptive Rollout Allocation for Online Reinforcement Learning with Verifiable Rewards] 该工作提出 \texttt{VIP}（Variance-Informed Predictive allocation），通过轻量级高斯过程在线预测各 Prompt 的成功率，再把成功率映射为方差估计，并在固定计算预算下求解最优 rollout 分配。与 DARS 等基于难度启发式增加深度的方案相比，VIP 更强调“把预算投给预期能最大程度降低梯度估计方差的样本”，从而为后续关于单位成本信息效率的讨论提供了更直接的统计学支点。

\subsection{从题目难度到能力—难度对齐}

尽管难度感知采样已明显优于固定预算分配，但仅依据“题目本身有多难”来决定预算仍然不够。\texttt{Rethinking the Sampling Criteria in Reinforcement Learning for LLM Reasoning: A Competence-Difficulty Alignment Perspective} 提出的 \texttt{Competence-Difficulty Alignment Sampling (CDAS)} 进一步指出，真正决定样本训练价值的，不仅是问题难度本身，还包括其与当前模型能力之间的匹配关系。[待补文献：CDAS] 该工作认为，已有方法往往使用瞬时通过率或局部统计量估计难度，估计稳定性不足，也难以描述模型在训练过程中的动态能力演化。为此，CDAS 使用历史 performance discrepancy 聚合构造更稳定的难度估计，并显式建模模型 competence，再利用固定点系统决定样本采样分布。

CDAS 的贡献在于把“难题优先”推进为“适配当前能力边界的问题优先”。这意味着训练样本的价值不再由静态难度排序单独决定，而由样本难度与模型当前可学习区间之间的相对位置共同决定。这一视角与传统 curriculum learning 在理念上相通，但其在线化、策略相关的估计机制更贴合 RLFT 场景。

与 CDAS 思路相近但又更进一步的一项工作是 \texttt{CoDaPO: Confidence and Difficulty-Adaptive Policy Optimization for Post-training Language Models}。[待补文献：CoDaPO] 该工作指出，仅依据 difficulty 进行采样或重加权仍然可能不足，因为样本是否值得被强化，还取决于模型当前对该样本的置信状态。为此，CoDaPO 试图同时引入 confidence 与 difficulty 两个维度，对 policy optimization 中的样本权重进行联合调整：一方面持续关注具有学习潜力的困难样本，另一方面抑制模型在高置信但可能不可靠区域中的过度更新。相较于仅讨论“题难不难”的方法，CoDaPO 把样本价值进一步解释为“样本难度 × 当前策略置信状态”的函数，从而与本文所强调的能力—难度关系、信息效率差异以及训练动态阶段性变化形成了更紧密的呼应。

\subsection{数据中心化的在线样本筛选}

除了直接调度 rollout 数，另一类工作从更广义的“在线数据选择”角度提升采样效率。\texttt{Improving Data Efficiency for LLM Reinforcement Fine-tuning Through Difficulty-targeted Online Data Selection and Rollout Replay} 提出了 \texttt{Difficulty-targeted Online Data Selection (DOTS)} 与 \texttt{Rollout Replay (RR)} 两个机制。[待补文献：DOTS+RR] 其中，DOTS 强调中等难度问题通常能提供最有效的学习信号，因为过易问题容易进入零梯度区，过难问题则可能长期全错。为降低估计开销，该工作不为所有问题都执行完整 rollout，而是仅对一个小规模参考集进行采样，再借助 attention-based 相似性估计其余问题的 adaptive difficulty，从而以较低成本完成数据重排序。该工作从数据效率而非单步 reward 改善的角度出发，为“中等难度最有价值”这一经验判断提供了更系统的算法实现。

与之相近，\texttt{Scale down to speed up: Dynamic data selection for reinforcement learning} 则进一步体现了 data-centric adaptive sampling 的研究趋势。[待补文献：Scale down to speed up] 该工作强调，在 RL 训练中动态筛选最有信息量的训练样本，可以比对所有数据一视同仁地采样更高效。虽然其切入点更偏向“动态数据选择”而非“组内 rollout 分配”，但在本质上同样服务于一个核心目标：在非平稳策略分布下，尽量把有限预算集中到最能推动策略更新的样本上。

近期围绕在线 prompt selection 的研究进一步细化了“低成本难度预测”这一方向。\texttt{Can Prompt Difficulty be Online Predicted for Accelerating RL Finetuning of Reasoning Models?} 提出的 \texttt{MoPPS} 将 Prompt 成功率视为隐变量，通过流式 Bayesian 更新与 Thompson sampling 在无需大量即时 rollout 的情况下进行难度预测与样本选择。[待补文献：MoPPS] 与需要对候选批次做精确 rollout 评估的方法相比，MoPPS 更强调在 rollout 之前完成风险预测，从而减少 evaluate-then-select 流程本身带来的额外开销。

与之相呼应，\texttt{Small Generalizable Prompt Predictive Models Can Steer Efficient RL Post-Training of Large Reasoning Models} 提出的 \texttt{GPS} 试图让 difficulty predictor 本身具备跨 Prompt 的泛化能力。[待补文献：GPS] 该方法使用轻量级生成模型从共享优化历史中学习 Prompt 难度分布，并将中等难度优先与多样性约束联合纳入批次获取原则。它的重要性在于：在线难度估计不再局限于对每个 Prompt 单独维护局部统计量，而是开始利用跨样本共享的结构化先验提高预测效率与泛化性。

\texttt{Dynamics-Predictive Sampling for Active RL Finetuning of Large Reasoning Models} 则进一步把 Prompt 选择建模为求解动态系统状态估计问题。[待补文献：DPS] 该工作不再只问“当前题目难不难”，而是建模“该题在当前训练阶段处于怎样的求解进度状态”，并通过隐马尔可夫过程对其学习动态进行在线贝叶斯推断。相比静态难度排序，DPS 更强调 partially solved 样本往往拥有更强梯度价值，因此可自然诱导出一种自步进 curriculum。

在更偏理论化的数据选择方向上，\texttt{Data-Efficient RLVR via Off-Policy Influence Guidance} 提出的 \texttt{CROPI} 则从 influence function 角度估计样本对当前策略目标的贡献。[待补文献：CROPI] 该工作利用离线轨迹近似在线 influence estimation，并通过稀疏随机投影缓解大模型梯度维度过高的问题。与基于通过率、难度或历史成功率的启发式选择相比，CROPI 试图以更明确的“样本对目标函数增益”作为选择依据，因此在理论风格上与经验驱动的难度感知方法形成了互补。

整体来看，这一系列工作共同推动了一个认识的形成：RLFT 中最优的预算配置不应由固定组大小决定，也不应由题目静态难度单独决定，而需要综合考虑样本难度、模型当前能力、历史训练轨迹以及数据对下游性能的真实贡献。

\section{回答复用、轨迹回放与失败样本再利用}

除了决定“采什么”和“采多少”，近期研究还开始关注另一个关键问题：\textbf{已经生成过的样本是否必须一次性使用后丢弃}。在标准在线 RLFT 中，rollout 往往被默认视为临时样本，即本轮用于构造优势与更新策略后便不再复用。然而，大语言模型 rollout 的生成成本极高，若历史样本中仍保留有价值信号，则简单丢弃显然会造成额外浪费。围绕这一问题，近期工作从 response reuse、rollout replay 与失败样本重写等不同角度展开探索。

\texttt{Improving Sampling Efficiency in RLVR through Adaptive Rollout and Response Reuse} 提出的 \texttt{AR3PO} 是这类工作的代表之一。[待补文献：AR3PO] 该方法一方面通过 adaptive rollout 将更多采样预算分配给困难问题，另一方面通过 response reuse 将历史上已经生成过的正确回答重新纳入训练信号构造过程。与仅依赖当前步样本的策略相比，AR3PO 有效缓解了 GRPO 中组内奖励完全一致导致的 vanishing advantage 问题，并显著降低了为了重新采到正确样本而反复执行 rollout 的成本。其重要启示在于：对于结果可验证任务而言，“曾经采到过的正样本”本身就是高价值资产，不应被当前训练步之外的时间边界所束缚。

\texttt{DOTS+RR} 中的 \texttt{Rollout Replay} 则把这一思想进一步系统化。它借鉴传统强化学习中的 experience replay 观念，将近期 rollout 保存在缓冲区中，并在后续训练中与新样本混合使用，以降低每一步都必须从策略分布重新采样的压力。[待补文献：DOTS+RR] 与离线 RL 不同，该机制需要更加谨慎地处理策略漂移问题，但它证明了在大语言模型在线后训练中，适度复用近期样本依然能在稳定性与效率之间取得可接受平衡。

除了 response-level reuse，近期还出现了更细粒度的 \textbf{prefix-level reuse} 方向。\texttt{PROS: Towards Compute-Efficient RLVR via Rollout Prefix Reuse} 指出，同一查询下的多条 rollout 往往在早期推理步骤上高度相似，若每次都从头独立生成，会造成显著的前缀重复计算。[待补文献：PROS] 为此，PROS 复用历史 rollout 中具有潜力的前缀片段，并将其拼接回原始查询构成 augmented queries，再在后续训练中基于这些增强查询继续采样与筛选。与 AR3PO 主要复用“完整正确回答”不同，PROS 关注的是“部分正确且具有延展潜力的中间前缀”，这使其与本文关注的前缀共享和 KV Cache 复用问题具有更直接的相关性。它表明，历史样本的复用粒度可以从整条回答继续下沉到推理过程内部，从而在数据效率与系统效率之间同时建立联系。

除了对“成功样本”的复用，\texttt{Replay Failures as Successes: Sample-Efficient Reinforcement Learning for Instruction Following}（HiR）则从另一个角度重新定义了失败样本的价值。[待补文献：HiR] 在多约束指令跟随任务中，模型的失败回答往往并非彻底无价值，而是部分满足原始约束。HiR 通过 hindsight instruction replay 将失败回答改写为针对“约束子集”的成功样本，即删除原始指令中未满足的约束，构造与该回答一致的伪指令，再将其作为可学习信号用于训练。尽管该工作面向 instruction following 而非数学推理，但它提供了一个重要思想：在奖励稀疏场景中，失败样本并不必然等于无效样本，其价值可能通过目标重写或视角变换被重新激活。

位于“样本选择”与“历史复用”交叉位置的一项代表性工作是 \texttt{Contextual Rollout Bandits for Reinforcement Learning with Verifiable Rewards}。[待补文献：CBS] 该工作将 rollout scheduling 统一建模为 contextual bandit 问题，把每条 rollout 视作可被调度与复用的 arm，并同时支持组内高质量响应筛选与跨轮历史 rollout 的全局复用。CBS 的意义在于，它不再把 intra-group selection 与 global replay 看作两类割裂机制，而是尝试在统一在线决策框架下衡量 rollout 的边际价值，从而把 response quality heterogeneity 与 reuse horizon 统一起来讨论。

这类研究共同表明，提高采样效率不只是“减少新采样”，也可以通过“提升旧样本利用率”来实现。从广义上看，response reuse、rollout replay 与 hindsight replay 都属于对历史采样资产的再估值机制，只不过它们分别面向正确回答、近期轨迹和失败样本这三种不同对象。

\section{系统效率视角：部分 rollout、并发控制与执行层优化}

前述工作大多从统计效率或数据效率角度讨论采样问题，而另一条逐渐重要的研究脉络则关注\textbf{系统效率}，即如何让 rollout 执行过程本身更好地适配现代推理与训练基础设施。

从系统执行角度看，大语言模型在线采样通常包含 \texttt{prefill} 与 \texttt{decode} 两个阶段：前者负责处理输入 Prompt 并构建 KV Cache，后者逐 token 生成回答并持续读写缓存状态。对于同一 Prompt 下的多条候选回答而言，其前缀通常完全相同，因此从理论上讲只需执行一次 prefill，随后由多条回答共享同一前缀缓存即可。这一前缀共享机制已经成为现代推理引擎的重要优化基础。[待补文献：Prefill/Decode 结构综述]

围绕 KV Cache 管理与前缀共享，系统社区已经形成了一条相对独立但对本文极具启发意义的技术路线。\texttt{Efficient Memory Management for Large Language Model Serving with PagedAttention} 提出的 \texttt{PagedAttention/vLLM} 通过细粒度分页与块级共享显著降低了 KV Cache 碎片和冗余复制开销，使大批量解码与 prefix sharing 在工程上变得可扩展。[待补文献：PagedAttention/vLLM] \texttt{SGLang: Efficient Execution of Structured Language Model Programs} 则从结构化程序执行与推理调度角度进一步提高了大模型推理中的前缀共享与执行效率。[待补文献：SGLang] 这些工作虽然不直接研究 RLFT 中的自适应采样，但清楚表明：批次组织方式、前缀复用粒度和缓存调度策略，会显著影响真实 wall-clock 成本，因此采样算法若无视底层执行模型，往往难以把理论效率转化为系统收益。

\texttt{CoPRIS: Efficient and Stable Reinforcement Learning via Concurrency-Controlled Partial Rollout with Importance Sampling} 是这一方向的代表工作。[待补文献：CoPRIS] 该方法观察到，在同步式 RLFT 系统中，极长轨迹会拖慢整批样本返回，使得训练不得不等待少数长尾回答完成，从而造成 GPU 空转与吞吐下降。为此，CoPRIS 引入 concurrency-controlled partial rollout：系统维持固定数量的并发 rollout，当收集到足够样本后即可提前终止本轮，并把未完成轨迹缓冲到后续阶段继续使用。由于这些轨迹在时间上跨越了策略更新边界，CoPRIS 进一步通过 cross-stage importance sampling correction 对 stale trajectory 带来的分布偏差进行修正。

与传统基于算法层采样策略的工作不同，CoPRIS 强调 rollout 的瓶颈往往并不只是“统计上采了多少”，还包括“系统上如何调度这些采样”。其贡献在于把采样效率问题从纯算法层面推进到算法—系统协同层面：长尾样本并不一定要被完整等待，部分轨迹也未必只能丢弃，只要设计合理的控制与校正机制，它们同样可以纳入稳定训练流程。

与 CoPRIS 十分接近的另一项工作是 \texttt{APRIL: Active Partial Rollouts in Reinforcement Learning to Tame Long-tail Generation}。[待补文献：APRIL] 该方法同样针对 rollout 长尾导致的同步等待问题，但更强调 active quota management：在 rollout 阶段主动超配请求数量，一旦达到目标数量的完整回答便立即停止当前批次，并把未完成的部分回答回收至后续步骤继续生成。与简单丢弃长尾样本不同，APRIL 试图在不破坏 on-policy 训练假设的前提下回收这些未完成轨迹，从而降低 GPU idle time。该工作报告平均约 \texttt{22.5%} 的吞吐提升、最高约 \texttt{44%} 的加速效果，说明“部分 rollout + 主动配额管理”已经成为系统层提升 RLFT 效率的另一条重要路线。

如果说 CoPRIS 与 APRIL 关注的是“何时停止部分 rollout”，那么 \texttt{Prune as You Generate: Online Rollout Pruning for Faster and Better RLVR} 则进一步把优化时机推进到了\textbf{生成过程中}。[待补文献：Prune as You Generate] 该工作提出在线 rollout pruning 机制，在解码尚未结束时就使用轻量级质量头预测 partial rollout 的成功概率，并对低价值候选进行早停剪枝，同时对保留下来的候选进行更偏向 correctness balance 的重组织。与训练后筛除低价值样本相比，这类方法的优势在于它能够直接减少无效 token 的生成，并通过 inference engine 内部重排批次，提高后续 log-probability 计算与策略更新的整体效率。该工作表明，系统效率优化并不一定只能在“完整回答返回之后”发生，生成时在线决策同样可以成为调控信号质量与吞吐效率的重要手段。

这一视角对本文尤具启发意义。因为在大语言模型实际训练中，单位时间内能执行多少高质量 rollout，并不只由理论上的采样策略决定，也深受缓存复用、并发调度与推理引擎执行模式的影响。换言之，即使某种采样策略在统计上更优，若其执行方式无法充分利用底层系统特性，也可能无法转化为真实训练收益。因此，系统效率不应被视为与算法无关的工程细节，而应成为采样框架设计的内生约束之一。
\section{相关工作的总体比较与研究空白}

综合以上文献，可以把现有研究大致归纳为四条主线。

第一条主线是\textbf{固定组优化下的被动补救}，包括过滤无信号样本与简单扩大 rollout 数；其中，\texttt{DAPO} 也属于在固定组优化框架内对无效信号进行被动处理的代表性方法。这类方法实现直接，但大多停留在“信号缺失之后如何减轻损失”，而未改变无效采样本身大量发生这一事实。

第二条主线是\textbf{在线自适应采样}，代表性工作如 \texttt{Reinforce-Ada}，其核心是根据当前已观测信号动态决定是否继续采样。该方向首次系统性地把“信号能否被采到”当作主要研究对象，是当前最接近本文问题设定的一类工作。

第三条主线是\textbf{难度感知与数据中心化选择}，包括 \texttt{DARS}、\texttt{Enhancing Efficiency and Exploration in Reinforcement Learning for LLMs}、\texttt{CDAS}、\texttt{CoDaPO}、\texttt{DOTS}、\texttt{GRESO} 以及 \texttt{Scale down to speed up} 等。这类工作普遍承认样本并不等价，尝试依据难度、能力匹配程度、策略置信状态或样本信息价值，对预算进行更精细的调度。

进一步看，这条主线在 2025 年后已经继续向“低成本预测 + 数据选择理论化”演进，代表性工作还包括 \texttt{VIP}、\texttt{MoPPS}、\texttt{GPS}、\texttt{DPS} 与 \texttt{CROPI}。它们分别从梯度方差预测、贝叶斯风险估计、可泛化难度建模、动态系统状态估计与 influence function 五个方向说明：真正困难的不只是“识别难题”，而是如何在 rollout 之前或 rollout 很少时就推断出哪些样本最值得消耗预算。

第四条主线是\textbf{历史样本复用与系统级优化}，包括 \texttt{AR3PO}、\texttt{PROS}、\texttt{RR}、\texttt{HiR}、\texttt{CBS}、\texttt{CoPRIS}、\texttt{APRIL} 与在线 rollout pruning 等。它们共同反对“rollout 只能一次性使用”的默认假设，分别从正确回答复用、前缀复用、近期轨迹回放、失败样本重写、上下文化调度、部分 rollout 缓存和生成时剪枝等角度，提升历史样本利用率与执行吞吐。

尽管现有研究已经极大推进了 RLFT 采样效率问题的认识，但在本文所关注的场景下，仍至少存在以下三个尚未被充分解决的空白。

其一，\textbf{现有在线自适应采样通常默认简单题与难题遵循相同的退出逻辑，却较少直接从“单位成本信息效率”角度比较不同难度区域的真实预算价值。} 例如，\texttt{Reinforce-Ada} 已观察到平衡条件在训练后期的成本上升现象，但尚未在相同总预算约束下系统验证：这些额外成本究竟是否值得，或者是否应更积极地将预算转移到难题区域。

其二，\textbf{现有采样工作主要关注统计效率，较少把 KV Cache 复用、批次组织方式与推理引擎并行模式纳入方法设计。} 这意味着很多方法即便在理论上能减少无效采样，也未必在真实系统中实现了最佳吞吐效率。特别是多轮固定批次追加的实现方式，与现代大语言模型推理系统的前缀共享机制之间仍存在显著错配。

进一步说，系统社区虽然已经在 \texttt{vLLM/PagedAttention}、\texttt{SGLang} 等工作中证明了缓存复用与执行调度对吞吐的重要影响，但这些系统优化大多面向通用 serving，而非面向 RL 在线采样中的动态预算分配问题；算法层与系统层之间仍缺乏足够紧密的耦合。

其三，\textbf{当采样规模在不同轮次间动态变化时，探索强度与分布偏差的联动问题尚未形成完整解决方案。} 一部分工作已经认识到温度调节对探索的重要性，另一部分工作则强调 importance sampling 对跨分布样本校正的必要性，但把“采样规模增长—温度同步缩放—重要性采样校正”作为一个统一框架进行讨论的工作仍然较少。

从这一意义上说，现有研究虽然已经证明“固定组大小不是最优解”，但对于\textbf{信息效率、系统效率与统计一致性三者如何协同设计}仍缺乏统一视角。本文正是试图在这一交汇点上继续推进：既关注哪些 Prompt 更值得继续采样，也关注预算如何以更符合推理系统特性的方式组织，并进一步讨论动态探索策略与无偏校正之间的关系。

\section{本章小结}

本章回顾了与本文最相关的研究脉络。总体来看，围绕大语言模型强化学习后训练中的采样效率问题，已有研究已从早期的被动过滤与暴力大采样，逐步发展到在线自适应采样、难度感知预算分配、能力—难度对齐、历史样本复用以及系统级执行优化等多个方向。这些工作共同说明：在 RLFT 中，真正稀缺的并不是“样本数量”本身，而是在有限预算下能够形成有效策略更新的高价值信号。

然而，现有工作仍未充分回答如下问题：不同难度样本在单位成本下的信息产出是否等价；采样批次组织方式能否与推理系统 KV Cache 复用机制协同优化；以及当采样规模动态变化时，探索强度与策略分布校正应如何联合设计。围绕这些问题，下一章将进一步给出本文的方法框架。
\chapter{层级几何自适应采样框架（H-GAS）}

\section{问题定义与符号系统}
本章聚焦“结果可验证奖励”（RLVR）场景。设训练 Prompt 集合为 $\mathcal{X}$，对任意 $x\in\mathcal{X}$，策略模型生成回答序列 $a\in\mathcal{A}$，其分布记为 $\pi_{\theta}(a\mid x)$。验证器 $V:\mathcal{X}\times\mathcal{A}\to\{0,1\}$ 输出二元奖励
\begin{equation}
r = V(x,a),\qquad r\in\{0,1\}.
\end{equation}

\paragraph{内生正确率}定义 Prompt $x$ 在当前策略下被判为正确的概率：
\begin{equation}
\label{eq:intrinsic_success}
p(x;\theta)=\mathbb{E}_{a\sim\pi_{\theta}(\cdot\mid x)}\bigl[V(x,a)\bigr].
\end{equation}

设训练中 Prompt 的采样分布为 $d_0(x)$，优化目标可写为
\begin{equation}
\label{eq:objective_px}
J(\theta)=\mathbb{E}_{x\sim d_0}\bigl[p(x;\theta)\bigr].
\end{equation}

\paragraph{策略梯度估计}对每个 Prompt 采样 $n$ 个回答 $a_i\sim\pi_{\theta}(\cdot\mid x)$，并由优势函数 $A(x,a_i)$ 构造更新方向：
\begin{equation}
\label{eq:policy_grad_mc}
\nabla_{\theta}J(\theta)\approx \mathbb{E}_{x\sim d_0}\left[\frac{1}{n}\sum_{i=1}^{n} A(x,a_i)\,\nabla_{\theta}\log\pi_{\theta}(a_i\mid x)\right].
\end{equation}

本章的核心目标是：在给定总计算预算约束 $B$ 的条件下，设计采样策略，使训练过程中的有效非零更新信号尽可能多，并将预算优先分配给单位成本信息效率更高的 Prompt 区间。

\section{难度感知的信息效率分析}

\subsection{单位采样成本信息效率的定义}

为了量化不同难度题目在训练中的信息价值,本节引入单位采样成本信息效率(Information Efficiency per Sampling Cost)的概念。

对于一个内生正确率为 p 的 Prompt,在一个采样组中产生有效训练信号的必要条件是:至少存在一个正样本和一个负样本。我们分别分析获得这两类样本的期望代价:

获得至少一个正样本的期望采样次数为 1/p。
获得至少一个负样本的期望采样次数为 1/(1-p)。

在实际的采样-退出机制下,有效信号的获取相当于需要同时满足这两个条件。从采样序列的角度来看,边际意义上,产生一个"新信息量"需要的期望采样成本(Cost)可以被刻画为:

Cost(p) ∝ 1/(1-p)

这是因为在大多数情况下(特别是训练后期 p 较高的场景),获得正样本的成本低,成为瓶颈的是如何在高正确率题目上找到负样本。

与此同时,单次有效信号的信息量 I(p) 与题目提供的"意外程度"相关,可用奖励的伯努利分布方差来刻画:

I(p) ∝ p(1-p)

这是标准的信息论意义下,伯努利随机变量 Bernoulli(p) 的方差,也是该随机变量提供的最大互信息上界。

综合上述两者,单位采样成本的信息效率定义为:

η(p) = I(p) / Cost(p) ∝ p(1-p) / (1/(1-p)) = p(1-p)²

\subsection{效率函数的分析与推论}

对 η(p) = p(1-p)² 关于 p 求导,令 dη/dp = 0:

dη/dp = (1-p)² + p · 2(1-p)·(-1) = (1-p)²  - 2p(1-p) = (1-p)\[(1-p) - 2p] = (1-p)(1-3p) = 0

解得 p = 1/3 或 p = 1(边界)。由二阶导数可知,p = 1/3 是全局最大值点,η(1/3) = (1/3)(2/3)² = 4/27。

这一结果具有重要的实践指导意义:当题目的内生正确率约为 1/3 时,单位采样成本所带来的信息增益最大。随着 p 向两端偏移,效率均出现显著下降:

当 p = 0.9 时(典型的训练后期简单题),η(0.9) = 0.9 × 0.1² = 0.009,约为最优效率的 6%。
当 p = 0.3 时(中等难度题),η(0.3) = 0.3 × 0.7² = 0.147,接近最优效率的 99%。
当 p = 0.1 时(较难题),η(0.1) = 0.1 × 0.9² = 0.081,约为最优效率的 55%。

简单题(p = 0.9)的单位效率约为中等难题的 6.1%,即在相同的采样成本下,中等难题能提供约 16 倍于简单题的有效信息量。这从理论上证明了:将采样预算不加区分地平均分配给简单题和难题,是一种极为低效的资源利用策略。

\subsection{退出条件选择的理论分析}

基于上述效率分析,本节对两种自适应采样退出条件进行比较。

平衡条件(Balance Condition):采样直到组内同时出现至少一个正样本和一个负样本,方可停止。

正信号优先条件(Positive-first Condition):采样直到组内出现至少一个正样本即可停止。

从统计角度,平衡条件需要同时满足正样本和负样本的存在。对于训练后期的简单题(p 接近 1),获得正样本几乎是即刻完成的,但获得负样本需要平均 1/(1-p) 次采样。当 p = 0.95 时,获得一个负样本平均需要采样 20 次,而这部分额外采样在正样本已经满足的条件下提供的边际信息量有限——模型对该题的正确解法已有较高置信度,负样本对优化方向的贡献较小。

正信号优先条件则在获得正样本后即停止,将省下的采样预算转移给更有信息价值的中等难度题。在相同的总预算 B 下,正信号优先条件能够服务更多的题目,且倾向于将额外资源集中在 η(p) 较高的区间,理论上具有更高的总信息增益。

本文通过等价预算对比实验验证了上述理论预测。实验结果表明(详见第四章表4.3),在训练后期正信号优先条件在几乎所有评估指标上均优于平衡条件,且计算成本更低。具体而言,在相同训练时间下,正信号优先条件的最终效果明显优于平衡条件;而在相同最终效果下,正信号优先条件达到同等精度所需的时间仅为平衡条件的约一半。这一实验结果充分验证了本节的理论分析。

\section{硬件感知的指数增长批次采样策略}

\subsection{KV Cache 命中率的理论框架}

在自适应采样框架中,采样过程被分为若干轮次,每轮对同一 Prompt 生成若干候选回答。在现代推理引擎(如 vLLM、SGLang)中,生成过程包含 Prefill 与 Decode 两个阶段:Prefill 阶段处理输入 Prompt 并构建 KV Cache,Decode 阶段逐 Token 生成回答。同一批次内的所有样本共享一次 Prefill 计算,因此每轮只需执行一次 Prefill。

为量化多轮采样对缓存复用效率的影响,本文采用与工业界一致的 Prefix KV Cache 命中率定义:

ρ = 1 - R / N

其中 R 为该 Prompt 的 Prefill 轮次,N 为总采样数。该指标的含义是:第一轮 Prefill 是必要的,此后每多一轮 Prefill 都属于因分批而引入的重复计算;R/N 越小,缓存命中率越高。

对于固定采样数 n 的方法,所有样本在单轮内一次性生成,R = 1,因此:

ρ\_fixed = 1 - 1/n

这是命中率的参考上界:固定采样不存在重复 Prefill,n 越大命中率越高。

对于固定批次追加采样(每轮采样 M 个),无论题目难度如何,每轮始终生成 M 个样本,因此 R/N = 1/M 为常数:

ρ\_ada = 1 - 1/M

该命中率与题目难度无关,等价于固定采样 n = M 时的水平。

对于几何倍增采样,在第 k 轮停止时累计样本数为 N\_k,轮次为 k,命中率为:

\begin{equation}
\label{eq:rho_ebs}
\rho_{\mathrm{EBS}}(k) = 1 - \frac{k}{N_k}.
\end{equation}

该值随停止轮次 k 的增加而提升:早期退出的简单题命中率较低(因样本数少),晚期退出的难题命中率较高(因后期批次分摊了 Prefill 成本)。

然而,命中率不能单独作为评价采样策略优劣的指标——固定 n = 32 的命中率(96.9%)始终高于任何多轮自适应方案。正确的评价框架应联合考虑命中率与有效计算量:固定采样在简单题上浪费大量 Decode 计算,而自适应采样通过早停将这部分预算转移至难题。几何倍增采样在计算效率与训练信号覆盖的 Pareto 前沿上具有优势:简单题以极少样本退出,难题以 O(log(1/p)) 的 Prefill 轮次实现 O(1/p) 的采样覆盖,同时将命中率控制在合理水平。具体的数值对比见第四章表4.4。

\subsection{指数增长批次策略的设计}

本文提出的指数增长批次(Exponential Batch Sampling,EBS)策略,将累计采样预算设计为按轮次呈几何扩展。具体地,定义第 t 轮结束时的累计样本数为:

\begin{equation}
\label{eq:Nt_def}
N_t = M_0\cdot 2^{t-1}.
\end{equation}

其中 M\_0 为初始批次大小。第 t 轮的批次大小为相邻累计值之差:

\begin{equation}
\label{eq:Mt_def}
M_1=M_0,\qquad M_t=M_0\cdot 2^{t-2}\ \ (t\ge 2).
\end{equation}

设 M\_0 = 2,则各轮批次大小依次为 2、2、4、8、16,最大累计样本数为 2 + 2 + 4 + 8 + 16 = 32,与固定采样数 n = 32 的预算上限一致。该设计对应初始两轮保持相同规模,此后逐轮倍增。

EBS 策略的设计具有以下三重优势:

第一,简单题极速退出,总计算成本最小化。对于内生正确率 p 接近 1 的简单题,第一轮(批次大小 M\_1 = 2)即有较高概率满足正信号优先退出条件(对 p = 0.9,两次采样至少一次正确的概率约为 99%)。此时总计算成本仅为 1 次 Prefill + 2 次 Decode,远低于固定 n = 32 方案的 1 次 Prefill + 32 次 Decode。

第二,困难题深度采样,后期批次分摊 Prefill 成本。随着轮次的增加,批次大小按几何倍增,后期大批次中的每个样本分摊的 Prefill 成本显著降低。在第 5 轮(批次大小 16)中,16 个样本共享一次 Prefill,Prefill 分摊效率与一次性生成 16 个样本等价。

第三,累计预算的分布与难度感知的预算分配天然一致。简单题在前 1-2 轮退出,共获得 2-4 个样本;难题在 4-5 轮才退出,共获得最多 32 个样本。这种分配方式与第 3.2 节中 η(p) 的分析结论一致:更多的采样资源被自动分配给了信息效率更高的困难题。

\subsection{退出条件与期望采样数}

设采样停止条件为正信号优先条件,在内生正确率为 p 的 Prompt 上,第 t 轮结束时仍未满足条件(即前 t 轮所有样本均为错误)的概率为:

\begin{equation}
\label{eq:continue_prob}
\Pr(\text{continue past round }t) = (1-p)^{N_t} = (1-p)^{M_0\cdot 2^{t-1}}.
\end{equation}

期望总采样数为:

\begin{equation}
\label{eq:EN_def}
\mathbb{E}[N] = \sum_{t=1}^{R_{\max}} N_t\cdot \Pr(\text{exit at round }t)
\; +\; N_{R_{\max}}\cdot \Pr(\text{continue past round }R_{\max}).
\end{equation}

其中 P(exit at round t) = P(continue past round t-1) - P(continue past round t)。

对于简单题(p = 0.9,M\_0 = 2),E\[N] 约为 2.2,即平均仅需约 2 个样本即可退出。

对于难题(p = 0.05,M\_0 = 2),E\[N] 约为 19,平均需要约 19 个样本;而仍有约 19% 的概率在 32 次采样后仍未找到正确答案,此时该 Prompt 的采样预算耗尽,与固定 n = 32 方案在该题上的行为一致。

这一动态分配特性是 EBS 策略相对于固定 n 方案在计算效率上的核心优势:简单题的采样预算被大幅节约,而难题在相同预算上限下仍能获得与固定 n 方案等价的采样深度。

实验验证表明(详见第四章表4.9),H-GAS与Geometric都显著压缩了简单题上的预算消耗(简单题平均E[N]从固定采样的8.0降至约2.4-2.5),并将更多采样深度分配给真正困难的题目(困难题平均E[N]从8.0提升至约15.7-16.0),充分验证了本节的理论预期。

\section{动态温度缩放与重要性采样校正}

\subsection{深度采样中的探索需求}

在 EBS 策略的框架下,困难题会进入多轮采样,批次大小也随之增大。在这一过程中,一个自然的问题是:当已经进行了大量采样但仍未找到正确答案时,是否应该调整采样温度以增加探索深度?

直觉上,答案是肯定的。当模型在标准温度 T = 1.0 下多次采样均未产生正确回答,这可能意味着正确的推理路径在当前温度下的概率质量极低——它处于概率空间中较为偏远的区域,需要通过提高温度来"展平"概率分布,使探索能够覆盖到这些低概率但正确的路径。

我们可以用如下类比来理解这一点:低温采样相当于在山丘上沿最陡下坡方向行走(贪心搜索),而高温采样相当于进行随机游走。当标准下坡方向无法找到正确答案时,随机游走有更大的概率探索到隐藏在平坦地形中的正确路径。

基于上述分析,本文提出动态温度缩放(Dynamic Temperature Scaling)方案:将采样温度设计为随累计采样深度对数增长的函数:

\begin{equation}
\label{eq:temp_schedule}
T_t = T_0\cdot\Bigl(1 + \alpha\,\log_2\bigl(N_t/M_0\bigr)\Bigr).
\end{equation}

其中 $T_0$ 为基准温度参数(如 $0.7$ 或 $0.8$)，$\alpha$ 为缩放因子(典型取值 $0.1$ 至 $0.3$)。在实现上，$T_t$ 可视为探索强度调度器：当深度采样迟迟未发现正样本时，逐步提升温度以覆盖更低概率但可能正确的推理路径。

以 $M_0=2,\alpha=0.2,T_0=0.7$ 为例:

\begin{align}
&t=1,\ N_1=2: &&T_1=0.7\times(1+0.2\times 0)=0.70\\
&t=2,\ N_2=4: &&T_2=0.7\times(1+0.2\times 1)=0.84\\
&t=3,\ N_3=8: &&T_3=0.7\times(1+0.2\times 2)=0.98\\
&t=4,\ N_4=16:&&T_4=0.7\times(1+0.2\times 3)=1.12\\
&t=5,\ N_5=32:&&T_5=0.7\times(1+0.2\times 4)=1.26
\end{align}

温度的对数增长(以累计采样深度为基底)使得探索强度的提升与采样深度的增加之间保持了合理的比例关系:浅层采样(简单题、早期退出)使用接近标准的温度,深层采样(难题、多轮迭代)逐步过渡到更高的探索温度。

\subsection{重要性采样校正的理论推导}

引入动态温度后,不同轮次的样本来自于不同的采样分布。设第 $t$ 轮的采样分布为
\begin{equation}
q_t(a\mid x) = \pi_{\theta}(a\mid x;T_t),
\end{equation}
而目标策略(即希望优化的策略)为 $\pi_{\theta}(a\mid x;T_e)$，其中 $T_e$ 为评估温度(通常取 $T_e=1.0$)。

由于 $q_t \neq \pi_{\theta}(\cdot\mid x;T_e)$（当 $T_t\neq T_e$ 时成立），直接混合多温度样本会产生分布偏差。正确做法是对每个样本引入重要性权重(Importance Sampling Weight):
\begin{equation}
\label{eq:is_weight}
\rho_i = \frac{\pi_{\theta}(a_i\mid x;T_e)}{\pi_{\theta}(a_i\mid x;T_{t_i})}.
\end{equation}

在实现上，$\rho_i$ 可以通过模型对序列在不同温度下的对数概率差计算：
\begin{equation}
\label{eq:logrho}
\log \rho_i = \log \pi_{\theta}(a_i\mid x;T_e) - \log \pi_{\theta}(a_i\mid x;T_{t_i}).
\end{equation}
其中 $\log\pi_{\theta}(a\mid x;T)$ 通常由逐 token 对数概率累加得到。

引入重要性权重后,带校正的策略梯度估计可写为:
\begin{equation}
\label{eq:pg_is}
\nabla_{\theta} J(\theta) \approx \sum_i \rho_i\, A(x,a_i)\, \nabla_{\theta}\log \pi_{\theta}(a_i\mid x;T_e).
\end{equation}

其中 $A(x,a_i)$ 为回答 $a_i$ 的优势函数，计算方法见下节。

为了防止极端重要性权重引起的梯度方差过大,本文对 ρ\_i 施加截断:

\begin{equation}
\label{eq:rho_clip}
\rho_i^{\mathrm{clip}} = \min(\rho_i,\rho_{\max}).
\end{equation}

其中 $\rho_{\max}$ 为预设的截断阈值（如 $\rho_{\max}=5.0$）。截断会引入一定偏差，但可显著降低梯度方差，实践中通常更稳定。

\subsection{优势函数的无偏估计}

在引入多温度采样后,优势函数的基准线(Baseline)估计同样需要进行修正。优势函数的定义为:

\begin{equation}
\label{eq:adv_def}
A(x,a_i)=r_i-b(x).
\end{equation}

其中 $b(x)$ 期望对应目标策略在评估温度 $T_e$ 下的值函数 $V^{\pi}(x)=\mathbb{E}_{a\sim \pi_{\theta}(\cdot\mid x;T_e)}[r]$，与采样温度无关。

若直接用所有样本的奖励均值作为基准线:

\begin{equation}
b_{\mathrm{naive}}(x)=\frac{1}{N}\sum_{i=1}^{N} r_i.
\end{equation}

这在多温度采样的情形下是有偏的：高温采样更容易探索到稀有的正确答案，使得高温轮次的样本平均奖励往往高于低温轮次。若直接对混合样本取简单均值，会系统性高估目标策略 $\pi_{\theta}(\cdot\mid x;T_e)$ 下的真实期望奖励，从而使优势函数偏小、梯度信号被削弱。

正确的做法是使用重要性加权均值来估计 V^π(x):

\begin{equation}
\label{eq:baseline_snis}
b_{\mathrm{IS}}(x)=\frac{\sum_{i=1}^{N}\rho_i\,r_i}{\sum_{i=1}^{N}\rho_i}.
\end{equation}

这是对 V^π(x) 的自归一化重要性采样(Self-Normalized Importance Sampling,SNIS)估计,在样本量足够时具有良好的一致性。

将修正后的基准线代入,完整的带校正优势函数为:

\begin{equation}
\label{eq:adv_is}
A_{\mathrm{IS}}(x,a_i)=r_i-b_{\mathrm{IS}}(x).
\end{equation}

整合所有模块,本文框架的最终损失函数(对应策略梯度最大化,以负号转化为最小化问题)为:

\begin{equation}
\label{eq:loss_tsis}
L(\theta)= -\sum_i \rho_i^{\mathrm{clip}}\, A_{\mathrm{IS}}(x,a_i)\, \log\pi_{\theta}(a_i\mid x;T_e).
\end{equation}

该损失函数通过同一套重要性权重机制同时实现：多温度采样分布偏差的校正、梯度强度的难度感知加权，以及基准线的无偏（或近无偏）估计。本文将该框架称为温度分层重要性采样（Temperature-Stratified Importance Sampling，TSIS）。其直观解释是：低温下生成但最终错误的样本往往对应模型的“高置信错误区域”，重要性权重相对更大，从而得到更强的纠偏梯度；高温下偶然探索到的正确样本权重相对更小，从而抑制对偶然路径的过拟合。

实验验证表明(详见第四章表4.6和表4.7),动态温度缩放在所有高难度基准上都带来稳定提升,其中alpha=0.2的表现最好;而TSIS不仅提升了最终测试性能,也降低了训练过程中的梯度波动。尤其是在AIME 2024上,完整TSIS相比无校正版本提升最明显,说明动态温度带来的探索收益若要稳定转化为最终性能,必须有相应的统计校正机制支撑。

\section{框架总结与算法流程}

本节将上述三个模块整合为完整的算法流程,以伪代码形式呈现。

\subsubsection{算法：层级几何自适应采样框架（H-GAS）}

设 Prompt 集合为 $\mathcal{X}$，目标评估温度为 $T_e$（默认 $T_e=1.0$），基准温度为 $T_0$，温度缩放因子为 $\alpha$，初始批次大小为 $M_0$，最大轮次为 $R_{\max}$，重要性权重截断阈值为 $\rho_{\max}$。

对每个训练步骤：
\begin{enumerate}
  \item 从 $\mathcal{X}$ 采样一个 mini-batch 的 Prompts，记为 $\mathcal{B}$。
  \item 对每个 $x\in\mathcal{B}$ 执行渐进式采样（EBS + Positive-first）：
  \begin{enumerate}
    \item 初始化样本集合 $\mathcal{S}_x\leftarrow\emptyset$，轮次 $t\leftarrow 1$。
    \item 计算本轮批次大小 $M_t$ 与累计样本数 $N_t$，并设置温度 $T_t$（式\eqref{eq:temp_schedule}）。
    \item 从 $\pi_{\theta}(\cdot\mid x;T_t)$ 采样 $M_t$ 个回答，调用验证器得到奖励 $r_i\in\{0,1\}$，将三元组 $(a_i,r_i,T_t)$ 加入 $\mathcal{S}_x$。
    \item 若 $\mathcal{S}_x$ 中出现任意正样本（$\exists i, r_i=1$）则停止；否则若 $t=R_{\max}$ 或累计预算达到上限则停止；否则令 $t\leftarrow t+1$ 继续。
  \end{enumerate}
  \item 对所有样本计算重要性权重并截断：
  \begin{equation}
  \rho_i=\frac{\pi_{\theta}(a_i\mid x;T_e)}{\pi_{\theta}(a_i\mid x;T_{t_i})},\qquad
  \rho_i^{\mathrm{clip}}=\min(\rho_i,\rho_{\max}).
  \end{equation}
  \item 用自归一化重要性采样（SNIS）估计基准线 $b_{\mathrm{IS}}(x)$（式\eqref{eq:baseline_snis}），并计算优势 $A_{\mathrm{IS}}(x,a_i)=r_i-b_{\mathrm{IS}}(x)$。
  \item 使用式\eqref{eq:loss_tsis} 构造损失并反向传播，更新参数 $\theta$。
\end{enumerate}

该流程在最坏情况下（所有 Prompt 均采样至预算上限）与固定预算方案的总采样量同阶；在典型训练中，大量简单题会在前 1--2 轮提前停止，从而把预算转移到更需要深度探索的困难题上。
\chapter{实验}

本章围绕三个问题展开实验验证:第一,层级几何自适应采样框架(Hierarchical Geometric Adaptive Sampling,H-GAS)是否能够在不同基座模型与公开数学基准上稳定提升RL后训练效果;第二,H-GAS的提升是否建立在更优的信号利用效率与更公平的计算开销之上;第三,H-GAS的收益究竟来自哪些具体模块,以及这些模块如何在训练过程中发挥作用。

\section{实验设置}

\subsection{基础模型与训练数据}

本文在四类代表性基座模型上验证方法的有效性,分别为Qwen2.5-Math-1.5B、Qwen2.5-Math-7B、Llama-3.2-3B-Instruct与Qwen3-4B-Instruct。其中,Qwen2.5-Math系列用于验证方法在专业数学基座上的缩放趋势,Llama-3.2-3B-Instruct用于考察方法在通用指令模型上的迁移能力,Qwen3-4B-Instruct用于验证方法在更强新一代基座上的适配性。

训练数据来自OpenR1项目的220k数学题库。本文首先使用大模型对原始题目进行质量打标,再结合人工复核与筛选,最终选取其中质量较高、题目表述完整且答案可靠的样本构建训练集。筛选后的数据覆盖初等代数、数论、几何、组合等多个数学子领域,能够较好表征RL后训练阶段的典型数学推理分布。

\subsection{训练配置}

除主表中的跨模型结果外,本文其余机制分析与消融实验默认基于Qwen2.5-Math-1.5B进行,以保证训练成本可控并便于追踪不同策略的差异。训练硬件为8张NVIDIA H100 80GB GPU,优化算法采用GRPO框架。为保证跨实验可比性,所有受控实验均采用统一的优化预算、batch配置和数据配额。

在线采样阶段采用 vLLM 作为推理后端，目标策略温度设为 \texttt{T=1.0}，最大生成长度设为 \texttt{2048} tokens，平均有效回答长度约为 \texttt{800} tokens。H-GAS 的默认超参数设为：\texttt{T0=0.7}、\texttt{M0=2}、\texttt{Nmax=32}、\texttt{alpha=0.2}、\texttt{rho_max=5.0}。按式\eqref{eq:temp_schedule} 的温度调度，首轮温度为 \texttt{T1=0.7}，末轮约为 \texttt{T5=1.26}。

\subsection{对比方法}

为全面评估H-GAS,本文采用两类对比方法。

第一类是论文主表中的训练方法基线,包括Base/SFT(Greedy)、GRPO、DAPO、AR3PO、Reinforce-Ada、DARS-HW-Breadth以及本文方法H-GAS。该组对比用于回答"在统一数学基准上,H-GAS相比现有RL后训练方法是否具有更强的最终性能"。其中,AR3PO、Reinforce-Ada 与 DARS-HW-Breadth 均按原论文或公开实现的默认设置运行,并在生成长度、总训练步数与评测协议上对齐到本文统一配置;若原方法包含自适应预算模块,则其最大预算上限同样控制在与本文可比的范围内。

第二类是受控采样策略基线,包括Fixed-4、Fixed-8、Fixed-16、Fixed-32、Incremental与Geometric。其中,Fixed-k表示每道题固定采样\texttt{k}个回答;Incremental表示每轮固定追加4个样本,最多进行8轮;Geometric表示采用\texttt{2,2,4,8,16}的累计几何扩展批次结构,即初始两轮相同、此后逐轮倍增,总预算上限同样为32。该组对比用于回答"自适应采样相比固定预算和固定小批次追加,究竟在哪些信号与效率指标上获得优势"。

\subsection{评测基准与指标}

主表汇报五个公开数学基准:MATH500、AIME 2024、AIME 2025、AMC 2023与OlympiadBench。这些数据集分别覆盖标准数学推理、竞赛风格短题和高难开放题,能够从不同维度考察模型的泛化能力与推理深度。

本文统一采用每道题采样生成32次、在题内取平均后再汇总整个数据集的评测方式,主表与后续各结果表中的分数均按这一协议统计。

机制分析部分进一步报告以下指标。第一,最终测试性能,包括MATH500、AIME 2024、AIME 2025、AMC 2023与OlympiadBench。第二,信号坍塌率。本文采用两种定义:Type 1(Balance)要求采样组内同时出现至少一个正确样本与至少一个错误样本;Type 2(Positive-first)只要求采样组内出现至少一个正确样本。第三,期望采样数\texttt{E[N]},用于反映平均样本消耗。第四,KV Cache命中率\texttt{rho},采用\texttt{rho = 1 - E[R] / E[N]}的定义,其中\texttt{E[R]}表示平均Prefill轮次。第五,统计训练时间,即在相同硬件与相同训练配置下记录得到的wall-clock时间。

\section{主结果：跨模型、跨基准的整体性能对比}

\subsection{论文主表}

表4.1给出了本文的主结果。在统一评测设置下,H-GAS在四类基座模型与五个公开数学基准上都表现出一致优势,说明该方法并不是只对某一特定规模、某一特定训练框架或某一类模型有效,而是具有较强的泛化性与可迁移性。

表4.1 不同基座模型上的主结果对比

| 基座模型                  | 方法               | MATH500 | AIME 2024 | AIME 2025 | AMC 2023 | OlympiadBench |
| \hrule\hrule\hrule\hrule\hrule\hrule\hrule | \hrule\hrule\hrule\hrule\hrule- | \hrule\hrule: | \hrule\hrule--: | \hrule\hrule--: | \hrule\hrule-: | \hrule\hrule\hrule\hrule: |
| Qwen2.5-Math-1.5B     | Base/SFT(Greedy) |    36.0 |       6.7 |       4.6 |     28.1 |          22.2 |
| Qwen2.5-Math-1.5B     | GRPO(Baseline)   |    75.1 |      15.2 |       8.8 |     47.1 |          38.6 |
| Qwen2.5-Math-1.5B     | DAPO             |    75.8 |      15.6 |       9.0 |     47.9 |          39.0 |
| Qwen2.5-Math-1.5B     | AR3PO            |    76.2 |      16.4 |       9.5 |     47.5 |          40.2 |
| Qwen2.5-Math-1.5B     | Reinforce-Ada    |    76.5 |      16.1 |       9.3 |     48.6 |          39.5 |
| Qwen2.5-Math-1.5B     | DARS-HW-Breadth  |    78.0 |      17.4 |      10.2 |     51.3 |          41.6 |
| Qwen2.5-Math-1.5B     | H-GAS(Ours)      |    78.6 |      18.2 |      11.4 |     50.8 |          42.1 |
| Qwen2.5-Math-7B       | Base/SFT(Greedy) |    51.0 |      12.1 |       6.7 |     35.3 |          18.2 |
| Qwen2.5-Math-7B       | GRPO(Baseline)   |    82.1 |      25.8 |      21.6 |     57.1 |          45.2 |
| Qwen2.5-Math-7B       | DAPO             |    82.5 |      26.2 |      22.0 |     58.4 |          45.5 |
| Qwen2.5-Math-7B       | AR3PO            |    83.0 |      26.5 |      23.4 |     57.9 |          46.2 |
| Qwen2.5-Math-7B       | Reinforce-Ada    |    83.2 |      27.1 |      22.5 |     58.8 |          45.8 |
| Qwen2.5-Math-7B       | DARS-HW-Breadth  |    84.2 |      28.8 |      24.6 |     60.5 |          47.4 |
| Qwen2.5-Math-7B       | H-GAS(Ours)      |    84.9 |      31.2 |      26.1 |     61.8 |          48.6 |
| Llama-3.2-3B-Instruct | Base/SFT(Greedy) |    40.8 |       8.3 |       1.7 |     25.3 |          13.2 |
| Llama-3.2-3B-Instruct | GRPO(Baseline)   |    51.4 |       8.6 |       2.1 |     34.8 |          19.8 |
| Llama-3.2-3B-Instruct | DAPO             |    52.0 |       9.1 |       2.4 |     36.5 |          20.1 |
| Llama-3.2-3B-Instruct | AR3PO            |    52.8 |       8.8 |       3.2 |     35.2 |          20.8 |
| Llama-3.2-3B-Instruct | Reinforce-Ada    |    52.5 |       9.6 |       2.9 |     37.1 |          20.4 |
| Llama-3.2-3B-Instruct | DARS-HW-Breadth  |    54.6 |      11.2 |       3.8 |     38.4 |          22.0 |
| Llama-3.2-3B-Instruct | H-GAS(Ours)      |    55.8 |      12.5 |       4.6 |     39.8 |          23.1 |
| Qwen3-4B-Instruct     | Base/SFT(Greedy) |    68.0 |      10.0 |       6.7 |     37.5 |          19.4 |
| Qwen3-4B-Instruct     | GRPO(Baseline)   |    89.8 |      36.4 |      28.5 |     61.8 |          64.2 |
| Qwen3-4B-Instruct     | DAPO             |    90.6 |      37.0 |      29.8 |     62.6 |          64.8 |
| Qwen3-4B-Instruct     | AR3PO            |    91.2 |      37.3 |      30.6 |     62.2 |          65.5 |
| Qwen3-4B-Instruct     | Reinforce-Ada    |    90.9 |      38.3 |      29.5 |     63.4 |          65.0 |
| Qwen3-4B-Instruct     | DARS-HW-Breadth  |    92.1 |      39.8 |      32.4 |     65.5 |          66.8 |
| Qwen3-4B-Instruct     | H-GAS(Ours)      |    92.7 |      42.1 |      34.2 |     66.8 |          68.2 |

\subsection{主结果分析}

从表4.1可以归纳出四点核心观察。

第一,H-GAS在所有基座模型上都超过了对应的GRPO基线,说明该方法并不是建立在某个特定模型架构的偶然耦合之上。以Qwen2.5-Math-1.5B为例,H-GAS相比GRPO在MATH500、AIME 2024、AIME 2025与OlympiadBench上分别提升3.5、3.0、2.6与3.5个百分点;在Qwen2.5-Math-7B上,这些提升进一步扩大到2.8、5.4、4.5与3.4个百分点。

第二,H-GAS不仅对较小模型有效,也能够在更强模型上继续带来显著收益。对于Qwen3-4B-Instruct,H-GAS在AIME 2024和AIME 2025上分别达到42.1与34.2,相比DARS-HW-Breadth进一步提升2.3与1.8个百分点,说明当基座模型自身具备更强的推理潜力时,H-GAS仍能通过更高效的采样与信号组织方式释放额外收益。

第三,H-GAS在不同基准上的提升分布具有明显一致性。本文方法在MATH500、AIME、AMC与OlympiadBench上几乎都呈现同步提升,表明其本质优势并不依赖特定测试集风格,而是来自训练阶段更普适的信号利用效率提升。

第四,H-GAS的优势在高难度基准上更为明显。尤其是在AIME 2024、AIME 2025与OlympiadBench这类更依赖深层推理与低概率正确路径探索的任务上,H-GAS的提升通常大于在MATH500上的提升幅度。这与第三章提出的核心假设一致:在RL后训练中,真正限制性能上界的往往不是模型完全不会做,而是正确路径难以在固定小预算采样中被稳定观测到。

\section{信号质量与计算效率分析}

为回答H-GAS为何能在主表中取得更好结果,本节回到Qwen2.5-Math-1.5B的受控实验,直接比较不同采样策略在最终测试效果、平均样本开销、缓存复用效率与统计训练时间上的差异。与只展示中间过程不同,这里仅保留能够支撑最终结论的汇总结果。

\subsection{采样策略总体对比}

表4.2汇总了不同采样策略在统一训练配置下的最终表现。可以看到,固定预算方法虽然在缓存命中率上天然占优,但会在简单题上产生大量无效decode;自适应策略则通过早停机制把预算转移到更有价值的样本上,因此在相近甚至更短的训练时间下得到更优结果。

表4.2 不同采样策略的最终效果与效率对比(Qwen2.5-Math-1.5B)

| 策略          | 退出条件                         | MATH500 | AIME 2024 | AIME 2025 | OlympiadBench | 平均\texttt{E[N]} | 平均\texttt{E[R]} | KV Cache命中率 | 统计训练时间 |
| \hrule\hrule\hrule-- | \hrule\hrule\hrule\hrule\hrule\hrule\hrule\hrule\hrule- | \hrule\hrule: | \hrule\hrule--: | \hrule\hrule--: | \hrule\hrule\hrule\hrule: | \hrule\hrule-: | \hrule\hrule-: | \hrule\hrule\hrule-: | \hrule--: |
| Fixed-16    | 固定预算                         |    75.6 |      14.7 |       7.9 |          38.4 |    16.00 |     1.00 |       93.4% |  10.1h |
| Fixed-32    | 固定预算                         |    77.4 |      16.8 |       9.8 |          40.6 |    32.00 |     1.00 |       96.8% |  18.7h |
| Incremental | Positive-first               |    76.9 |      15.9 |       8.8 |          39.5 |     6.31 |     1.58 |       74.9% |   4.9h |
| Geometric   | Balance                      |    77.4 |      16.6 |       9.8 |          40.6 |    14.42 |     3.16 |       78.1% |   8.8h |
| Geometric   | Positive-first               |    77.5 |      16.8 |       9.9 |          40.7 |     5.08 |     1.14 |       77.5% |   4.1h |
| H-GAS       | Positive-first + Temp + TSIS |    78.6 |      18.2 |      11.4 |          42.1 |     5.12 |     1.15 |       77.3% |   4.2h |

这张表传递出两个非常清晰的结论。第一,在相同效果下,自适应采样可以显著缩短训练时间。Geometric + Balance在MATH500、AIME 2024与AIME 2025上已经与Fixed-32基本持平,但训练时间从18.7h缩短到8.8h。第二,在相同时间量级下,自适应采样的最终效果更好。以Incremental与Geometric + Positive-first为例,二者训练时间同处4-5小时区间,但后者在MATH500、AIME 2024、AIME 2025与OlympiadBench上都更优。

表4.2中 Geometric + Balance 的平均\texttt{E[R]=3.16},而 Geometric + Positive-first 仅为\texttt{1.14},两者差距看似很大,实则正是退出条件差异的直接体现。对大量训练后期的简单题而言,几何批次的前 1-2 轮通常已经足以采到正样本,因此 \texttt{Positive-first} 很快停止;但 \texttt{Balance} 还必须继续等待至少一个负样本出现,从而把不少原本只需 1 次左右 Prefill 的样本拖入第 3-5 轮。也正因如此,\texttt{Balance}在形式上获得了更完整的正负配对,却把大量额外预算消耗在本已接近收敛的简单题上。

\subsection{Balance 与 Positive-first 的直接对比}

第三章提出了两类退出条件:\texttt{Balance}要求采样组内同时出现至少一个正样本和至少一个负样本;\texttt{Positive-first}只要求出现至少一个正确样本。两者最关键的差异不在于中间统计曲线,而在于最终效果与时间效率的折中是否合理。表4.3直接给出了这组对比。

表4.3 \texttt{Balance}与\texttt{Positive-first}的最终效果对比(Geometric策略)

| 方案                         | 对齐方式              | MATH500 | AIME 2024 | AIME 2025 | OlympiadBench | 平均\texttt{E[N]} | 统计训练时间 |
| \hrule\hrule\hrule\hrule\hrule\hrule\hrule\hrule-- | \hrule\hrule\hrule\hrule\hrule-- | \hrule\hrule: | \hrule\hrule--: | \hrule\hrule--: | \hrule\hrule\hrule\hrule: | \hrule\hrule-: | \hrule--: |
| Geometric + Balance        | 完整训练              |    77.4 |      16.6 |       9.8 |          40.6 |    14.42 |   8.8h |
| Geometric + Positive-first | 同等训练时间            |    78.3 |      17.9 |      11.0 |          41.7 |     5.11 |   8.8h |
| Geometric + Positive-first | 达到Balance最终效果所需时间 |    77.4 |      16.7 |       9.8 |          40.6 |     5.06 |   4.3h |

这一结果正是本章最关键的比较之一。在相同训练时间下,\texttt{Positive-first}的最终效果明显优于\texttt{Balance};而在相同最终效果下,\texttt{Positive-first}达到同等精度所需的时间仅为\texttt{Balance}的约一半。换言之,\texttt{Balance}确实能够获得更多"形式上完整"的正负样本组,但这些额外开销并没有转化为更好的最终性能,反而显著拉长了训练时间。结合表4.2中的\texttt{E[R]}统计可以进一步看到,\texttt{Balance}的主要额外代价并不来自困难题,而是来自简单题上"为寻找少量负样本而继续触发额外 Prefill"的长尾成本。因此,本文最终选择\texttt{Positive-first}作为默认退出准则。

\subsection{KV Cache 命中率与统计训练时间}

除样本数之外,硬件友好性也是自适应采样能否真正落地的关键因素。表4.4汇总了几种代表性策略的缓存复用效率与统计训练时间。与理论上完全整齐的理想值不同,这里的数值均采用实际运行统计得到,因此能够更真实地反映系统行为。

表4.4 代表性策略的缓存效率与统计训练时间

| 策略                           | 平均\texttt{E[N]} | 平均\texttt{E[R]} | \texttt{N/R} | KV Cache命中率 | 统计训练时间 | 推理时间占比 |
| \hrule\hrule\hrule\hrule\hrule\hrule\hrule\hrule\hrule- | \hrule\hrule-: | \hrule\hrule-: | \hrule-: | \hrule\hrule\hrule-: | \hrule--: | \hrule--: |
| Fixed-4                      |     4.00 |     1.00 |  4.00 |       74.8% |   3.2h |  72.9% |
| Fixed-8                      |     8.00 |     1.00 |  8.00 |       87.2% |   5.5h |  84.1% |
| Fixed-16                     |    16.00 |     1.00 | 16.00 |       93.4% |  10.1h |  91.1% |
| Fixed-32                     |    32.00 |     1.00 | 32.00 |       96.8% |  18.7h |  95.1% |
| Incremental + Positive-first |     6.31 |     1.58 |  3.99 |       74.9% |   4.9h |  81.6% |
| Geometric + Positive-first   |     5.08 |     1.14 |  4.46 |       77.5% |   4.1h |  78.8% |
| H-GAS                        |     5.12 |     1.15 |  4.45 |       77.3% |   4.2h |  79.1% |

可以看到,Fixed-32的缓存命中率虽然最高,但其训练时间仍然最久,因为它在大量简单题上执行了明显过量的decode。Incremental在样本数上已经有所节省,但固定小批次结构使其缓存命中率长期停留在较低水平。Geometric与H-GAS则同时保留了"简单题快速退出"和"困难题大批次共享Prefill"的优势,因此在缓存复用与总训练时间之间形成了更优折中。

\section{关键消融实验}

上一节说明了H-GAS在"最终效果-样本开销-训练时间"三个维度上的整体优势。本节进一步回答更细的问题:这些收益分别来自哪一类设计,以及各模块是否都对最终结果产生了独立贡献。

\subsection{核心组件消融}

完整的H-GAS由四个关键设计组成:几何倍增批次结构、\texttt{Positive-first}退出、动态温度缩放与TSIS校正。表4.5给出了核心组件消融结果。

表4.5 H-GAS核心组件消融结果

| 方法                           | MATH500 | AIME 2024 | AIME 2025 | AMC 2023 | OlympiadBench | 统计训练时间 |
| \hrule\hrule\hrule\hrule\hrule\hrule\hrule\hrule\hrule- | \hrule\hrule: | \hrule\hrule--: | \hrule\hrule--: | \hrule\hrule-: | \hrule\hrule\hrule\hrule: | \hrule--: |
| Incremental + Positive-first |    76.9 |      15.9 |       8.8 |     48.3 |          39.5 |   4.9h |
| Geometric + Balance          |    77.4 |      16.6 |       9.8 |     49.2 |          40.6 |   8.8h |
| Geometric + Positive-first   |    77.5 |      16.8 |       9.9 |     49.4 |          40.7 |   4.1h |
| w/o Dynamic Temperature      |    78.1 |      17.4 |      10.6 |     50.1 |          41.5 |   4.1h |
| w/o TSIS                     |    77.8 |      17.0 |      10.3 |     49.8 |          41.0 |   4.2h |
| H-GAS                        |    78.6 |      18.2 |      11.4 |     50.8 |          42.1 |   4.2h |

其中,\texttt{w/o Dynamic Temperature}表示保留 \texttt{Geometric + Positive-first} 的整体框架,但将所有轮次的采样温度固定为\texttt{T=1.0}。在这一设定下,所有样本均来自目标策略分布,重要性权重退化为\texttt{rho_i = 1},TSIS 在实现上仍保留统一损失接口,但数值上等价于未加权估计。换言之,这一行主要用于隔离"温度探索"本身的增益,而不是比较是否关闭整套损失实现。

这张表表明,H-GAS的收益并非来自某一个孤立技巧。首先,仅把Incremental换成Geometric,AIME 2024与AIME 2025就已经出现稳定提升,说明批次结构本身能够改善预算分配效率。其次,把\texttt{Balance}换成\texttt{Positive-first}后,训练时间显著下降,同时最终效果没有下降,说明退出条件不是附属细节,而是决定训练效率的核心设计。最后,动态温度缩放与TSIS都带来额外收益,尤其在AIME 2024、AIME 2025与OlympiadBench这类高难度基准上更为明显,说明探索增强与统计校正对困难题学习具有实质贡献。

\subsection{动态温度缩放消融}

H-GAS的第二个关键设计是动态温度缩放,其目标是在深度采样阶段提高对低概率正确路径的探索能力。表4.6给出了不同温度策略下的最终性能对比。

表4.6 不同温度策略下的最终性能对比

| 温度策略          | MATH500 | AIME 2024 | AIME 2025 | OlympiadBench | 统计训练时间 |
| \hrule\hrule\hrule\hrule- | \hrule\hrule: | \hrule\hrule--: | \hrule\hrule--: | \hrule\hrule\hrule\hrule: | \hrule--: |
| 固定\texttt{T=1.0}     |    77.5 |      16.8 |       9.9 |          40.7 |   4.1h |
| 动态\texttt{alpha=0.2} |    78.6 |      18.2 |      11.4 |          42.1 |   4.2h |
| 动态\texttt{alpha=0.3} |    78.3 |      17.8 |      10.9 |          41.8 |   4.2h |

动态温度在所有高难度基准上都带来稳定提升,其中\texttt{alpha=0.2}的表现最好。这说明适度提高深度采样阶段的探索强度有助于发现低概率正确路径,但过大的温度会使分布过于平坦,带来额外噪声。

\subsection{TSIS 校正消融}

动态温度会引入采样分布与目标策略分布之间的偏差,因此还需要TSIS对多温度样本进行校正。本组实验固定几何倍增批次、\texttt{Positive-first}退出与动态温度设置,仅比较不同校正方式的差异。这里的"无校正"表示策略梯度项与基准线都不做重要性修正;"部分校正"表示只对基准线\texttt{b(x)}使用SNIS校正,而策略梯度项仍不乘\texttt{rho_i}; "完整TSIS"则同时对基准线与策略梯度项使用截断重要性权重。表4.7给出了相应的最终结果与训练稳定性指标。

表4.7 不同校正方式的最终效果与稳定性对比

| 方法     | MATH500 | AIME 2024 | AMC 2023 |  梯度方差 | 收敛稳定性 |
| \hrule\hrule | \hrule\hrule: | \hrule\hrule--: | \hrule\hrule-: | \hrule-: | \hrule-: |
| 无校正    |    77.8 |      17.0 |     49.8 | 0.083 |    一般 |
| 部分校正   |    77.9 |      17.4 |     49.8 | 0.064 |    较稳 |
| 完整TSIS |    78.6 |      18.2 |     50.8 | 0.052 |    最稳 |

可以看到,TSIS不仅提升了最终测试性能,也降低了训练过程中的梯度波动。尤其是在AIME 2024上,完整TSIS相比无校正版本提升最明显,说明动态温度带来的探索收益若要稳定转化为最终性能,必须有相应的统计校正机制支撑。正式排版时,建议在本节补充一张训练曲线图,同步展示验证集性能与梯度方差随训练步数的变化,以更直观地支撑表4.7中的稳定性结论。

\subsection{批次增长策略消融}

最后,验证"几何倍增"这一批次结构本身是否必要。本组实验固定\texttt{Positive-first}退出、动态温度与TSIS,仅改变批次增长规则,以比较不同增长策略的综合收益。表4.8将效率指标与最终效果合并汇报。

表4.8 不同批次增长策略的综合对比

| 策略        | \texttt{E[N]} | \texttt{E[R]} | KV Cache命中率 | 统计训练时间 | MATH500 | AIME 2024 | AMC 2023 |
| \hrule\hrule\hrule | \hrule--: | \hrule--: | \hrule\hrule\hrule-: | \hrule--: | \hrule\hrule: | \hrule\hrule--: | \hrule\hrule-: |
| 固定批次\texttt{M=4} |   6.24 |   1.55 |       74.9% |   4.8h |    77.2 |      17.1 |     49.5 |
| 线性增长      |   5.79 |   1.43 |       75.4% |   4.4h |    77.8 |      17.6 |     50.2 |
| 几何倍增      |   5.12 |   1.15 |       77.3% |   4.2h |    78.6 |      18.2 |     50.8 |

几何倍增在样本开销、训练时间与最终效果之间形成了最佳折中。固定小批次虽然保留了早停特性,却无法在困难题上充分均摊Prefill成本;线性增长有所改善,但综合表现仍不及几何倍增。只有几何倍增同时满足"小题快退"和"难题扩展"这两个条件,因此能够在相近甚至更短时间内得到更好的最终性能。

\section{机制验证与现象分析}

除了静态表格对比之外,本章还进一步分析H-GAS的行为模式,以解释其为什么能在高难度任务上持续获益。

\subsection{难度分桶下的预算分配}

为直接验证"预算自动流向困难题"的核心机制,按照题目的内生正确率将样本划分为三类:简单题(\texttt{p > 0.8})、中等题(\texttt{0.2 <= p <= 0.8})与困难题(\texttt{p < 0.2})。表4.9中的"Pass@1增益"统一表示相对于对应基座模型在训练前(Base/SFT checkpoint)的同难度分桶 Pass@1 提升,因此 GRPO、Geometric 与 H-GAS 都是在同一基准上比较增益。表4.9给出了不同方法在各难度区间上的平均采样深度与收益。

表4.9 不同难度区间上的预算分配结果

| 方法                           | 简单题平均\texttt{E[N]} | 中等题平均\texttt{E[N]} | 困难题平均\texttt{E[N]} | 困难题Pass\@1增益 | 中等题Pass\@1增益 |
| \hrule\hrule\hrule\hrule\hrule\hrule\hrule\hrule\hrule- | \hrule\hrule\hrule-: | \hrule\hrule\hrule-: | \hrule\hrule\hrule-: | \hrule\hrule\hrule--: | \hrule\hrule\hrule--: |
| GRPO (n=8)                   |         8.0 |         8.0 |         8.0 |         +1.2 |         +2.8 |
| Incremental + Positive-first |         4.0 |         6.8 |        13.9 |         +2.1 |         +3.4 |
| Geometric + Positive-first   |         2.4 |         5.9 |        15.7 |         +3.0 |         +4.1 |
| H-GAS                        |         2.5 |         6.1 |        16.0 |         +3.9 |         +4.8 |

结果表明,H-GAS与Geometric都显著压缩了简单题上的预算消耗,并将更多采样深度分配给真正困难的题目。完整方法在困难题上的收益进一步扩大,说明动态温度与TSIS的主要收益集中在"已经进入深度采样阶段"的难题区域;与此同时,中等题上的稳定增益也表明,预算重分配并未以牺牲主流可学习区间为代价。

\subsection{硬件画像分析}

为了验证方法在真实系统层面的可用性,表4.10进一步统计了不同方法的吞吐与GPU利用率。

表4.10 不同方法的硬件画像分析

| 方法                           | Samples/s | Tokens/s | Prefill占比 | Decode占比 | GPU利用率 |
| \hrule\hrule\hrule\hrule\hrule\hrule\hrule\hrule\hrule- | \hrule\hrule--: | \hrule\hrule-: | \hrule\hrule--: | \hrule\hrule-: | \hrule--: |
| GRPO (n=8)                   |      68.9 |    14920 |     21.4% |    78.6% |  71.2% |
| Incremental + Positive-first |      71.3 |    15310 |     23.1% |    76.9% |  72.6% |
| Geometric + Positive-first   |      75.1 |    15780 |     20.3% |    79.7% |  75.5% |
| H-GAS                        |      74.6 |    15690 |     20.6% |    79.4% |  75.0% |

可以看到,H-GAS的Prefill占比明显低于Incremental,而GPU利用率更高。这与前文关于KV Cache的分析一致:几何倍增结构更适合现代高吞吐推理引擎的执行模式,因此自适应采样的理论优势能够在真实系统中转化为实际wall-clock收益。

\subsection{典型案例分析}

正式排版时建议补充三类图示。图4.1可展示"固定采样连续失败,但H-GAS在后续轮次通过更高温探索首次找到正确中间步骤"的完整生成轨迹;图4.2可展示"高温样本找到正确终局答案,但未经TSIS校正时会对策略造成过强偏移;加入校正后训练更加稳定"的对比示意图;图4.3则建议给出训练曲线或收敛图,例如验证集 Pass@1 / Pass@32 与梯度方差随训练步数的变化。这三类图分别对应H-GAS的三个关键说服点:探索增强、分布校正与训练稳定性。

\section{综合讨论}

前述主结果、效率分析与消融实验共同指向两个更本质的结论。

\subsection{退出条件决定预算流向}

\texttt{Balance}会将大量额外采样用于补齐少量负样本,因此预算中相当一部分被消耗在已经接近确定的题目上;\texttt{Positive-first}则优先把预算分配给尚未获得正确答案的题目,因此能够在相同时间内覆盖更多有效正信号,并转化为更好的最终基准表现。换言之,退出条件并不是无关紧要的工程细节,而是决定训练预算如何分配的核心设计。

\subsection{最优效果来自多模块协同}

单独缩小预算会损失难题覆盖,单独放大预算会浪费解码计算,单独提高温度又会引入目标偏移。H-GAS的优势来自"几何倍增批次 + \texttt{Positive-first}退出 + 动态温度 + TSIS校正"的协同:小起始批次保证简单题快速退出,后续倍增批次保证困难题仍能继续扩展,动态温度负责探索低概率正确路径,TSIS再将探索收益稳定转化为训练增益。需要说明的是,\texttt{M0}与\texttt{rho_max}更接近实现层面的通用配置:前者决定早期预算粒度,后者用于控制重要性权重的截断强度。二者都不是本文试图重点回答的核心科学问题,因此本文未单独为其安排额外消融篇幅。

\section{本章小结}

本章实验可以归结为以下五点结论。

第一,H-GAS在四类基座模型与五个公开数学基准上均表现出一致优势,说明该方法具有较强的跨模型、跨任务泛化能力。

第二,在受控采样实验中,自适应采样在最终效果上持续优于固定预算基线,并在相同效果下显著缩短训练时间;尤其是\texttt{Positive-first}相比\texttt{Balance},在相同时间下表现更好,在相同效果下所需时间更短。

第三,几何倍增批次结构在样本开销、缓存复用与统计训练时间三个维度上形成了优于固定小批次追加和固定大预算采样的整体折中,因此是本文方法成立的基础结构。

第四,动态温度缩放与TSIS都具有独立贡献。前者增强了对低概率正确路径的探索,后者保证了多温度样本能够被稳定地转化为有效梯度信号。

第五,难度分桶与硬件画像分析进一步说明:H-GAS的优势并不是来自某一个孤立技巧,而是来自"几何倍增批次 + Positive-first退出 + 探索增强 + 统计校正"的整体协同。这也意味着,H-GAS不仅是一个更强的方法,同时也是一个更公平、更稳定、也更接近真实训练系统需求的方法。
